%% 第31回編集委員会（物理第5巻）：組 原W7
%% p5-14.tex  原-14　放射線をはかる・身を守る：計測と防護（★1〜4・3段にしない）
%% 中身は計画書 tools/phys5-plan.md の B「原-14」。公式の持ち主は A（A7・A13・A70・A76・A77）。
%%   ほかの章は計画書 B の節の番号で指した（原-9 ⑦、原-12 ④⑨、原-13 ②③⑤⑩⑪、原-15、原-16、波-7 ②、熱-1 ③）。
%% 数値はすべて Python で計算した：checks/物理第5巻/原-14/calc_p5_14.py（出力 calc_p5_14_out.txt）、
%%   常用漢字の確かめ kanji_p4_13.py（joyo.txt は2010年の2136字）。
%% 前の巻・数学の本との約束（計画書 A-6・C-6-2）
%%   ・M3（6-4⑨ ガイガーカウンター・ポアソン分布）→ ③（ph:原-14:sec3）
%%   ・M14（3-14⑤・3-17：放射線を数えると σ≈√N、1分 100±10、10分 1000±32）→ ③⑨（同じ数を使った）
%%   ・M18（8-8⑬ 立体角、S17）→ ⑥・例題11
%%   ・A76（波-7 ② の逆2乗の法則）→ ⑥
%%   ・熱-1 ③ の比熱 → 例題2
%% 記号（この章）：吸収線量 D（Gy）、等価線量 H_T（Sv）、実効線量 E（Sv）、放射線加重係数 w_R、組織加重係数 w_T、
%%   放射能 A（Bq）、距離 r、計数率 m（はかった）・n（本当の）、不感時間 τ、数えた数 N、立体角 Ω、検出器の半径 a、
%%   線量率定数 Γ。実効線量 E とぶつからないよう、エネルギーは E_γ（γ線）、E_α のように添字で書く（N-4）。
%%   原-13 の減弱係数 μ をそのまま使う。原-13 の断面積 σ とは章が別で、この章の σ は標準偏差だけ。
%% 出典（11月11日に確かめた）
%%   ・環境省「放射線による健康影響等に関する統一的な基礎資料（令和7年度版）」HTML 版
%%       （https://www.env.go.jp/chemi/rhm/current/kisoshiryohtml.html の各ページを curl で読んだ。表はスライドの画像で読んだ）
%%       01-01-03（放射線は五感で感じないが測定は容易）、01-03-09・01-03-10（透過力と体の中での影響の範囲）、
%%       02-03-03〜02-03-08（Bq・Gy・Sv、w_R：γ・X・β 1、陽子 2、α・重イオン 20、中性子 2.5〜21、w_T：骨髄（赤色）・結腸・
%%       肺・胃・乳房 0.12、生殖腺 0.08、膀胱・食道・肝臓・甲状腺 0.04、骨表面・脳・唾液腺・皮膚 0.01、残りの組織の合計 0.12、
%%       ICRP 2007年勧告。物理量・防護量・実用量、実用量は安全側に少し大きく定義、1 cm 線量当量）、
%%       02-04-01〜02-04-06（測定器：GM 計数管・電離箱・NaI(Tl) シンチレーション・Ge 半導体・個人線量計、検出限界、3σ 法
%%       （こえる確率 約0.1 %）、バックグラウンド 0.1 μSv/h 程度、地上 1 m ではかる）、02-04-10（低減係数：木造家屋 約0.4）、
%%       02-04-13〜02-04-15（預託実効線量：成人50年・子ども70歳まで。経口の預託実効線量係数 成人：Sr-90 0.028、I-131 0.022、
%%       Cs-134 0.019、Cs-137 0.013、Pu-239 0.25、トリチウム 0.000018 μSv/Bq：ICRP Publ. 119）、
%%       02-05-01〜02-05-04（自然 2.1 mSv/年（日本）・2.4（世界）、医療 2.6、合計 4.7。内訳 宇宙線 0.3、大地 0.33、
%%       ラドン222 0.37、ラドン220 0.09、食品の鉛210・ポロニウム210 0.80、カリウム40 0.18、炭素14 0.014 ほか：
%%       原子力安全研究協会「生活環境放射線 第3版 増補版」2024年）、02-05-02（生活空間の空間線量率 0.01〜1 μSv/h）、
%%       02-05-11（カリウムの 0.01 % が放射性）、02-05-15（早見図）、02-05-17（トリチウム 12.3年・生物学的半減期 約10日）、
%%       03-01-02〜03-01-04（確定的影響と確率的影響、しきい線量は1 %の人に症状が出る線量）、03-03-04（しきい値の表）、
%%       03-04-03（100 mSv あたり 0.5 %、1000人中 約300人ががんで亡くなる）、04-01-08（LNT をめぐる論争：米国科学アカデミー
%%       BEIR VII 2006・NCRP 2018 と、フランスの医学・科学アカデミー 2005）、04-01-09・04-01-13・04-02-04（線量限度は
%%       安全と危険の境界ではない、一般公衆 1 mSv/年）、04-02-03（食品の基準値）、04-03-01（外部被ばくの低減三原則）。
%%   ・電離放射線障害防止規則 第3条（管理区域：3月 1.3 mSv）・第4条（5年 100 mSv・1年 50 mSv・女性 3月 5 mSv）・
%%       第5条（眼の水晶体・皮膚）・第7条（緊急作業 100 mSv）：e-Gov 法令 API
%%       https://laws.e-gov.go.jp/api/2/law_data/347M50002000041 の本文を読んだ。
%%   ・空気の質量エネルギー吸収係数 μen/ρ：NIST "X-Ray Mass Attenuation Coefficients" ComTab/air.html
%%       （0.6 MeV 0.02953、0.8 MeV 0.02882 cm²/g）。
%%   ・Cs-137 の 1 cm 線量当量率定数 0.0927 μSv·m²/(MBq·h)：日本アイソトープ協会『アイソトープ手帳』の値。原典の表は
%%       作業環境から届かず、二次資料2つ（Isotope News 2012年5月号の記事、量研の放射線診療 Q&A：どちらも検索の抜粋）で照らし合わせた。
%%       NIST の μen/ρ から計算した空気カーマ率定数 0.076 μGy·m²/(MBq·h) の約1.2倍で、実用量が安全側に大きく定義される
%%       ことと合う（calc_p5_14.py）。
%%   ・Cs-137（Ba-137m）の γ 線 0.662 MeV・放出の割合 約85 %：原-13 の頭のコメントと同じ二次資料。
%% 〔仮〕例題のために置いた値（本文に「とします」）：GM 計数管の不感時間 100 μs、トリチウム水 1500 Bq/L、線源の強さ、
%%   検出器の大きさ、背景の計数率、照度計の値。
%% 〔要確認〕高校の「物理基礎」で Gy・Sv を扱うか（学習指導要領の本文は「放射線の種類と性質、放射性物質の基本的な性質」
%%   までで、単位の扱いは解説・教科書で確かめられなかった）。U-5 により「この章で導く公式」の欄には入れなかった。
\makeatletter
\tikzset{txo lab/.style={font=\sffamily\footnotesize, text=bkInk}, txo arr/.style={line width=1pt, -{Stealth[length=2.4mm]}},
  txo dim/.style={bkGray, line width=0.5pt, {Stealth[length=1.8mm]}-{Stealth[length=1.8mm]}},
  txo box/.style={draw=bkInk, line width=0.5pt, rounded corners=1mm, inner sep=2.5pt, font=\sffamily\footnotesize, align=center}}
\makeatother

\phchapter{14}{放射線をはかる・身を守る：\\計測と防護}{目に見えない放射線を、どうやってはかる？}{\bklevelrange{1}{4}}{放射線をはかる・身を守る：計測と防護}
\bkmichishirube{3-11,3-14,3-17,6-4,8-8,波-7,熱-1,原-9,原-12,原-13}{原-15,原-16}

\begin{goal}
\begin{itemize}[label={\color{bkTheme}●}, leftmargin=1.6\zw, itemsep=0pt, topsep=1pt]
\item \bklevel{1}\enspace 放射線は目に見えないが、電離を使って数えられることと、Bq・Gy・Sv の意味のちがいを説明できる。
\item \bklevel{2}\enspace 吸収線量から等価線量・実効線量を計算し、自然放射線や医療の線量と比べられる。
\item \bklevel{3}\enspace 点線源の逆2乗の法則と遮へいを組み合わせ、時間・距離・遮へいの3原則で線量を見積もれる。
\item \bklevel{4}\enspace 計数のばらつき（$\sqrt N$）・数え落とし・検出限界・立体角を使って、測定の結果を正しく読める。
\item \bklevel{2}\enspace 確定的影響と確率的影響のちがい、線量限度の意味、低い線量のリスクについての2つの見方を説明できる。
\end{itemize}
\end{goal}

\phkoushiki{%
\phk{放射線は目に見えないが、物を電離する性質を使ってはかれる（放射線の性質と利用）}{中学理科}{ph:原-14:sec2}
}

\begin{hatena}
\noindent\begin{minipage}[c]{0.52\linewidth}\raggedright
放射線は、目に見えず、音もせず、においもありません。体にあたっても、何も感じません。

それなのに、ニュースや自治体のサイトには「\qty{0.05}{\micro Sv/h}」のような数字がならび、病院や空港では、放射線の量をきちんとはかって管理しています。感じることもできない物を、どうやって数えるのでしょう。そして、はかった数字は、何を表しているのでしょう。
\end{minipage}\hfill
\begin{minipage}[c]{0.45\linewidth}\centering
\begin{tikzpicture}[x=1mm, y=1mm]
  \filldraw[fill=bkGray!20, draw=bkInk, line width=0.6pt, rounded corners=2mm] (4,4) rectangle (40,30);
  \filldraw[fill=white, draw=bkInk, line width=0.5pt] (8,16) rectangle (36,26);
  \node[font=\sffamily\bfseries\small, text=bkInk] at (22,21) {0.05};
  \node[txo lab] at (30,8.5) {\si{\micro Sv/h}};
  \filldraw[fill=bkGray!45, draw=bkInk, line width=0.5pt] (40,12) rectangle (54,20);
  \node[txo lab, anchor=north] at (47,11) {検出部};
  \foreach \y in {8,16,24} {\draw[bkOrange!85!black, line width=0.7pt, decorate, decoration={snake, amplitude=0.7mm, segment length=2.2mm, post length=1.6mm}, -{Stealth[length=1.6mm]}] (58,\y+6) -- (54.5,\y/3+12);}
  \node[txo lab, anchor=south] at (22,31) {測定器（サーベイメーター）};
\end{tikzpicture}
\end{minipage}
\end{hatena}

\begin{yattemiyou}[距離を2倍にすると、明るさは？（光で放射線のまね）]
\textbf{用意するもの}\par
小さな LED のライト（光る所が小さいもの）、照度計のアプリを入れたスマホ、メジャー\par\vspace{2pt}
\noindent\begin{minipage}[c]{0.60\linewidth}
\begin{enumerate}[label={\bkcircled[bkGreen]{\arabic*}}, leftmargin=2.4\zw, itemsep=1pt, topsep=0pt]
\item 部屋を暗くし、スマホを机に置く。LED を、照度計のセンサーの真上 \qty{20}{cm} に、下向きに固定して明るさ（lx）を読む。
\item 距離を \qty{40}{cm}・\qty{60}{cm}・\qty{80}{cm} と変えて読む。ライトを消したときの明るさも読み、それを引く。
\item 自治体や国が公開している、空間線量率（\si{\micro Sv/h}）の地図やモニタリングポストの値を、家族といっしょに、住んでいる所と、ほかの地方とで見くらべる。
\end{enumerate}
\end{minipage}\hfill
\begin{minipage}[c]{0.37\linewidth}\centering
\begin{tikzpicture}[x=1mm, y=1mm]
  \filldraw[fill=bkGray!25, draw=bkInk, line width=0.5pt] (14,34) rectangle (22,38);
  \node[txo lab, anchor=west] at (23,36) {LED};
  \foreach \a in {-25,-12,0,12,25} \draw[bkYellow!90!black, line width=0.6pt, -{Stealth[length=1.5mm]}] (18,33.5) -- ++(\a-90:22);
  \filldraw[fill=bkInk!75, draw=bkInk, line width=0.5pt, rounded corners=1mm] (6,4) rectangle (30,9.5);
  \node[font=\sffamily\footnotesize, text=white] at (18,6.8) {照度計};
  \draw[txo dim] (34,9.5) -- (34,34);
  \node[txo lab, anchor=west] at (34.5,22) {$r$};
\end{tikzpicture}
\end{minipage}
\par\vspace{2pt}\textbf{気をつけること}：放射線の測定器を使うときは、学校や科学館で、先生や係の人といっしょに使います（①）。LED の光を、目に向けないようにしましょう。
\end{yattemiyou}

\begin{bkcenter}
\begin{bktable}{colspec={Q[l,m,wd=34mm] Q[c,m,wd=14mm] Q[c,m,wd=14mm] Q[c,m,wd=14mm] Q[c,m,wd=14mm]}, rowsep=1pt, row{1}={bg=bkTheme!10}, cells={font=\small}}
距離 $r$〔cm〕 & 20 & 40 & 60 & 80 \\
明るさ〔lx〕 & 1610 & 395 & 180 & 102 \\
20 cm のときの何分の1か & 1 & 4.1 & 8.9 & 15.8 \\
\end{bktable}
\bktabcaption{やってみようで、ミナさんが調べた結果（ライトを消したときの明るさを引いた値）}\label{ph:原-14:tab-yattemiyou}
\end{bkcenter}

\noindent やってみようの答え：表\ref{ph:原-14:tab-yattemiyou}から、距離を2倍にすると明るさは約 $\tfrac14$、3倍で約 $\tfrac19$、4倍で約 $\tfrac1{16}$ になっています。明るさは\textbf{距離の2乗に反比例}します。小さな所から四方八方に出た光が、距離 $r$ の所では面積 $4\pi r^2$ の球面全体に広がるからです。音の強さ（物理シリーズの『波動』（波-7 ②））と同じ\textbf{逆2乗の法則}で、小さな放射線の出どころ（点線源）から出る放射線も、まったく同じように弱まります（⑥）。距離をとることが、放射線から身を守るいちばん簡単な方法の1つなのです（⑦）。

前の章（原-13）では、放射線が物にあたると、通り道の原子を電離しながらエネルギーを渡すことを学びました。この章では、その電離を使って放射線を\textbf{数え}、体が受けたエネルギーを\textbf{線量}として表し、線量を\textbf{減らす}方法を考えます。放射線について大切なのは、「あるか、ないか」「危ないか、安全か」の二択で考えずに、「\textbf{どれだけの量か}」で考えることです。私たちは、自然からも、毎日ある量の放射線を受けています。その量と比べて、ほかの量がどれくらいかを、数で読めるようになりましょう。

\begin{omoidasou}
\bkrecall{原-13}{放射線は、通り道の原子から電子をはじき出す（電離）。空気で1組のイオンを作るのに使うエネルギーは約 34 eV。γ線は $e^{-\mu x}$ で減り、半価層は $\ln2/\mu$。}
\bkrecall{原-12}{1秒に1個の原子核がこわれる量が \qty{1}{Bq}（ベクレル）。放射性の原子は、半減期ごとに半分になる。}
\bkrecall{波-7}{点の音源から出る音の強さは、距離の2乗に反比例する（逆2乗の法則）。}
\bkrecall{3-14}{数えた数 $N$ がたまにしか起きないできごとの数なら、ばらつき（標準偏差）は約 $\sqrt N$。1分で100なら $\pm10$、10分で1000なら約 $\pm32$。}
\bkrecall{6-4}{ガイガーカウンターが鳴る回数は、ポアソン分布にしたがう（⑨）。}
\bkrecall{8-8}{半分の角が $\alpha$ の円すいの立体角は $\Omega=2\pi(1-\cos\alpha)$。まわり全体は $4\pi$。}
\bkrecall{熱-1}{水 \qty{1}{kg} の温度を \qty{1}{K} 上げるには \qty{4.18e3}{J} の熱が要る（比熱）。}
\end{omoidasou}

\noindent この章の記号：この章では、$E$ を\textbf{実効線量}（⑤）に使います。そのため、γ線のエネルギーは $E_\gamma$ のように添字をつけて書き、ふつうのエネルギーは「エネルギー」とことばで書きます。$\sigma$ は標準偏差（3-14）で、原-13 の断面積ではありません。$\mu$ は原-13 と同じ減弱係数です。

%% ======================== 物理の道具 ========================
\bksection{doguu}{物理の道具}

\bksubsection{1}{はかる前に：安全のための約束}
\label{ph:原-14:sec1}
\begin{column}{安全のために：放射線をはかるときの約束}
\begin{itemize}[label={\color{bkTheme}●}, leftmargin=1.6\zw, itemsep=1pt, topsep=1pt]
\item 放射線を出す物（鉱物・古い時計の光る文字盤・古いガス灯の部品など）を、自分で集めたり、こわしたり、体に近づけたりしない。
\item 放射線の測定器は、学校・自治体・科学館が持っている物を、先生や係の人といっしょに使う。立ち入りが禁止されている場所や、「管理区域」の標識がある場所には入らない。
\item X線の装置、霧箱（ドライアイスや液体の窒素を使う）、高い電圧を使う管の実験は、自分ではしない。科学館や病院で、専門の人が見せてくれるときに見る。
\item はかった数字だけで「危ない」「安全」と決めつけない。単位（\si{\micro Sv/h}・mSv・Bq）を確かめ、自然放射線などの量と比べて考える（⑩）。
\end{itemize}
\end{column}
この約束は、こわがるためのものではありません。放射線の影響は「受けた量」で決まります（⑪）。必要のない量をむやみに受けないこと、そして、数字を正しく読むことが、身を守るための第一歩です。

\Needspace{8\baselineskip}\bksubsection{2}{放射線を数えるしくみ：検出器}
\label{ph:原-14:sec2}
放射線をはかる道具（\textbf{検出器}）は、どれも、放射線が通り道で起こす\textbf{電離}（原-13 ②）か、電子を高いエネルギーの状態にする\textbf{励起}を、電気の信号に変えます。\bkidx{けんしゅつき}{検出器}

\begin{bkfigure}
\begin{tikzpicture}[x=1mm, y=1mm]
  % GM 計数管
  \filldraw[fill=bkBlue!8, draw=bkInk, line width=0.8pt] (10,4) rectangle (70,24);
  \draw[bkInk, line width=1.2pt] (6,14) -- (74,14);
  \node[txo lab, anchor=north] at (40,2.5) {金属の筒（$-$）・中のガス・まん中の細い線（$+$）};
  \draw[bkOrange!85!black, line width=0.9pt, decorate, decoration={snake, amplitude=0.8mm, segment length=2.4mm, post length=2mm}, -{Stealth[length=2mm]}] (24,34) -- (32,6);
  \node[txo lab, anchor=south] at (24,34.5) {放射線};
  \node[font=\sffamily\scriptsize, text=bkRed!80!black] at (25.5,17.5) {$+$};
  \node[font=\sffamily\scriptsize, text=bkBlue!80!black] at (31.5,18) {$-$};
  \draw[bkBlue!80!black, line width=0.6pt, -{Stealth[length=1.5mm]}] (32.5,17.4) -- (35,15);
  \foreach \x in {38,40,42,44,46} {\node[font=\sffamily\scriptsize, text=bkBlue!80!black] at (\x,15.8) {$-$};}
  \foreach \x in {38,40,42,44,46} {\node[font=\sffamily\scriptsize, text=bkBlue!80!black] at (\x,12.2) {$-$};}
  \node[txo lab, anchor=west] at (48,19.5) {電子なだれ};
  % 回路
  \draw[bkInk, line width=0.6pt] (74,14) -- (86,14) -- (86,24);
  \draw[bkInk, line width=0.6pt] (86,14) -- (102,14);
  \filldraw[fill=white, draw=bkInk, line width=0.6pt] (78,24) rectangle (100,32);
  \node[txo lab] at (89,28) {高い電圧};
  \filldraw[fill=bkGray!25, draw=bkInk, line width=0.6pt] (102,9) rectangle (122,19);
  \node[txo lab] at (112,14) {数える};
  \node[txo lab, anchor=north] at (112,7.5) {1個 → 1つのパルス};
\end{tikzpicture}
\bkfigcaption{GM 計数管（ガイガー・ミュラー計数管）のしくみ（模式図）。放射線がガスを電離すると、電子は高い電圧で線に向かって加速され、ほかの分子を次々に電離して、雪だるま式に増える（電子なだれ）。放射線1個が、1つの大きな電気のパルスになる。}\label{ph:原-14:fig-gm}
\end{bkfigure}

\begin{bkcenter}
\begin{bktable}{colspec={Q[l,m,wd=23mm] Q[l,m,wd=33mm] Q[l,m,wd=33mm] Q[l,m,wd=27mm]}, rowsep=1pt, row{1}={bg=bkTheme!10}, cells={font=\small}}
検出器 & しくみ & 得意なこと & 使われる所 \\
電離箱 & ガスの電離でできたイオンと電子を、そのまま電流として集める & 強い放射線でも、数え落としなく線量率をはかる & 放射線を使う施設、線量計の基準 \\
GM 計数管 & ガスの電離を電子なだれで大きくし、1個ずつのパルスにする & 弱い放射線でも1個ずつ数えられる。β線に感度が高い & 体や物の表面の汚れを調べる \\
シンチレーション検出器 & 結晶（NaI など）が放射線で光り、その光を電気に変えて大きくする & γ線をよくとらえる。光の量からエネルギーがわかる & 空間線量率、食品の検査 \\
半導体検出器 & 半導体（ゲルマニウムなど）の中の電離でできた電子と正孔を集める & エネルギーを細かく見分けられる & 食品・土の放射能（どの核種か） \\
\end{bktable}
\bktabcaption{主な検出器（環境省の基礎資料 2.4節をもとにまとめた）。}\label{ph:原-14:tab-kenshutsu}
\end{bkcenter}

\phboxhead{\bkhatenaref の答え}
\label{ph:原-14:hatena}放射線は目に見えませんが、通り道の原子を\textbf{電離}します。ガスをつめた管に高い電圧をかけておくと、電離でできた電子が電極に集まり、とちゅうで電子なだれを起こして、\textbf{放射線1個が、1つの電気のパルス}になります。これを数えるのが GM 計数管です。光る結晶（シンチレーター）では、放射線のエネルギーが光に変わり、その光を電気に変えて数えます。半導体（原-9 の pn 接合に逆向きの電圧をかけたもの）では、電離でできた電子と正孔が、そのまま電流になります。

こうして数えた数（1秒あたりの数）から、体が受けるエネルギー（Gy）や、影響の大きさの目安（Sv）を計算します（④⑤）。測定器が \si{\micro Sv/h} を表示しているのは、数えた数に、決まった係数（校正定数）をかけて換算しているのです。だから、測定器は、使う前に基準の線源で正しく合わせておく（\textbf{校正}する）必要があります。

\phboxhead{パルスの大きさで、エネルギーがわかる}
GM 計数管は、どんなエネルギーの放射線でも、同じ大きさのパルスを出します（電子なだれが、いつも同じくらいまで育つため）。だから、数はわかっても、エネルギーはわかりません。シンチレーション検出器と半導体検出器のパルスの大きさは、放射線が落としたエネルギーに比例します。パルスの大きさごとに数を分けて並べると、核種ごとに決まったエネルギー（セシウム137 なら $E_\gamma=\qty{0.662}{MeV}$）の所に山ができ、「何の核種が、どれだけあるか」がわかります。半導体では、電子と正孔の1組を作るのに要るエネルギーが数 eV と、空気の約 34 eV（原-13 ②）よりずっと小さいので、同じエネルギーで多くの組ができ、山がするどくなります。

\Needspace{8\baselineskip}\bksubsection{3}{数えた数のばらつきと、数え落とし}
\label{ph:原-14:sec3}
放射線は、いつ来るかが偶然で決まります。同じ場所で1分ずつはかっても、数は毎回ちがいます。平均して $N$ 個数えられるとき、数のばらつき（標準偏差 $\sigma$）はおよそ
\pheq{\sigma\approx\sqrt{N}}{個 ＝ $\sqrt{\text{個}}$（数えた数）}
\label{ph:原-14:sqrtN}です（3-14・6-4⑨のポアソン分布）。1分で100個なら $100\pm10$（ばらつき 10 \%）、10分で1000個なら $1000\pm32$（3.2 \%）です。ばらつきの割合は $\sqrt N/N=1/\sqrt N$ なので、\textbf{長く数えるほど、割合は小さく}なります。

\begin{phbaisuru}[4倍に]{数える時間}
数える数 $N$ が4倍になり、ばらつき $\sqrt N$ は2倍になりますが、ばらつきの割合 $1/\sqrt N$ は $\tfrac12$ になります。割合を $\tfrac1{10}$ にしたければ、時間は100倍要ります。精しくはかるには、とても長い時間がかかるのです。
\end{phbaisuru}

\phboxhead{数え落とし：パルスのあとの「休み時間」}
GM 計数管は、1つのパルスを出したあと、管の中のイオンが片づくまでの短い時間、次の放射線を数えられません。この時間を\textbf{不感時間} $\tau$ といいます。\bkidx{ふかんじかん}{不感時間}1秒に $m$ 個数えたとすると、1秒のうち $m\tau$ 秒は「休み時間」で、数えられるのは残りの $1-m\tau$ 秒だけです。本当に来た放射線の数を1秒に $n$ 個とすると、数えた数は $m=n(1-m\tau)$ なので、
\pheq{n=\frac{m}{1-m\tau}}{1/s ＝ (1/s) ÷ 1}
\label{ph:原-14:fukan}です。

\begin{phbaisuru}{数えた数 $m$}
\noindent\begin{minipage}[c]{0.46\linewidth}\raggedright
$\tau=\qty{100}{\micro s}$ のとき、$m$ が1秒に50個なら数え落としは 0.5 \%、500個なら 5 \%、2000個なら 20 \% です。数えた数が2倍になると、数え落としの割合（$m\tau$）も2倍になります。本当の数がどれほど多くても、数えた数は1秒に $1/\tau=\num{10000}$ 個をこえません。強い放射線の場所で GM 計数管が少なめに表示するのはこのためで、そういう所では電離箱を使います。
\end{minipage}\hfill
\begin{minipage}[c]{0.50\linewidth}\centering
\begin{tikzpicture}
\begin{axis}[bkplot, axis lines=left, width=50mm, height=42mm, xmin=0, xmax=6000, ymin=0, ymax=12000, xtick={2000,4000}, ytick={4000,8000}, xlabel={$m$（/s）}, ylabel={$n$（/s）}]
  \addplot[bkcurve, domain=0:5500, samples=80] {x/(1-x*1e-4)};
  \addplot[bkGray, dashed, line width=0.8pt, domain=0:6000] {x};
  \draw[bkGray, dash pattern=on 1.5pt off 1pt, line width=0.5pt] (axis cs:2000,0) -- (axis cs:2000,2500) -- (axis cs:0,2500);
  \node[font=\sffamily\footnotesize, anchor=north west, inner sep=1pt] at (axis cs:4200,3000) {$n=m$（破線）};
\end{axis}
\end{tikzpicture}
\end{minipage}
\end{phbaisuru}

\Needspace{8\baselineskip}\bksubsection{4}{体が受けたエネルギー：吸収線量（Gy）}
\label{ph:原-14:sec4}
放射線に関する単位は、「出す側」と「受ける側」に分けられます。\textbf{出す側}の単位は、原-12 で学んだ放射能の \textbf{Bq}（1秒にこわれる原子核の数）です。\textbf{受ける側}の単位の1つめは、物 1 kg が放射線から受け取ったエネルギーです。
\pheq{D=\frac{\Delta E_{\text{吸}}}{\Delta m}}{Gy ＝ J/kg}
\label{ph:原-14:Gy}これを\textbf{吸収線量}といい、\bkidx{きゅうしゅうせんりょう}{吸収線量}単位は\textbf{グレイ}（Gy＝J/kg）です。$\Delta E_{\text{吸}}$ は、質量 $\Delta m$ の部分が吸収したエネルギーです。吸収線量は、熱量計などで直接はかることもできる、\textbf{物理量}です。

1 Gy は、どれくらいのエネルギーでしょう。体重 60 kg の人の全身が 1 Gy を受けると、受け取るエネルギーは $1\times60=\qty{60}{J}$ で、体温の上がり方は
\[ \frac{60}{60\times4.18\times10^3}\approx\qty{0.00024}{K} \]
にしかなりません（例題2）。熱としてはごくわずかです。ところが、1 Gy は、短い時間に受けると、血をつくる骨髄のはたらきが落ち始める量（しきい線量 約 0.5 Gy：⑪）をこえる量です。放射線の影響は、熱ではなく、\textbf{電離で DNA などの分子をこわす}ことによるからです。60 J は、約 34 eV ずつに分けると、体の中で約 $10^{19}$ 組もの電離を起こすエネルギーなのです。

\Needspace{8\baselineskip}\bksubsection{5}{影響の大きさで重みをつける：等価線量と実効線量（Sv）}
\label{ph:原-14:sec5}
同じ吸収線量でも、放射線の種類によって、体への影響の大きさがちがいます。α線は、原-13 ⑤で見たように、通り道にぎっしり電離を起こすので、同じエネルギーでも、細胞の DNA の近い所を何か所もまとめて傷つけやすく、影響が大きくなると考えられています。そこで、臓器や組織ごとの吸収線量 $D_{\mathrm T}$ に、\textbf{放射線加重係数} $w_{\mathrm R}$ をかけた量を、\textbf{等価線量}といいます。\bkidx{とうかせんりょう}{等価線量}
\pheq{H_{\mathrm T}=w_{\mathrm R}D_{\mathrm T}}{Sv ＝ 1 × Gy}
\label{ph:原-14:HT}単位は\textbf{シーベルト}（Sv）です。\bkidx{しーべると}{シーベルト}

さらに、臓器や組織ごとに、放射線でがんなどが起きやすいかどうかがちがいます。等価線量に\textbf{組織加重係数} $w_{\mathrm T}$ をかけて、全身について合計したものを、\textbf{実効線量}といいます。\bkidx{じっこうせんりょう}{実効線量}
\pheq{E=\textstyle\sum_{\mathrm T}w_{\mathrm T}H_{\mathrm T}}{Sv ＝ 1 × Sv}
\label{ph:原-14:E}$w_{\mathrm T}$ の合計はちょうど1になるように決められているので、全身が同じ等価線量を受けたときは、実効線量はその等価線量と同じ値になります。実効線量は、体の一部だけが受けたときも、体の外から受けたときも、体の中から受けたときも、同じ物差しで足し合わせられる量として作られたのです。

\begin{bkcenter}
\begin{bktable}{colspec={Q[l,m,wd=50mm] Q[c,m,wd=14mm] Q[l,m,wd=46mm] Q[c,m,wd=10mm]}, rowsep=1pt, row{1}={bg=bkTheme!10}, cells={font=\small}}
放射線 & $w_{\mathrm R}$ & 組織・臓器 & $w_{\mathrm T}$ \\
γ線・X線・β線 & 1 & 骨髄・結腸・肺・胃・乳房 & 0.12 \\
陽子線 & 2 & 生殖腺 & 0.08 \\
α線・重いイオン & 20 & ぼうこう・食道・肝臓・甲状腺 & 0.04 \\
中性子線 & 2.5〜21 & 骨の表面・脳・だ液腺・皮ふ & 0.01 \\
 & & 残りの組織の合計 & 0.12 \\
\end{bktable}
\bktabcaption{放射線加重係数と組織加重係数（国際放射線防護委員会 ICRP の2007年勧告。環境省の基礎資料による）。中性子線は、エネルギーによって変わる。}\label{ph:原-14:tab-w}
\end{bkcenter}

\Needspace{14\baselineskip}
\begin{phflow}[吸収したエネルギーから、実効線量へ]
\phflowrow{$D=\dfrac{\Delta E_{\text{吸}}}{\Delta m}$}{放射線の種類で\\重みをつける}{$H_{\mathrm T}=w_{\mathrm R}D_{\mathrm T}$}
\phflowrow{$H_{\mathrm T}=w_{\mathrm R}D_{\mathrm T}$}{組織の感じやすさで\\重みをつけて合計}{$E=\textstyle\sum w_{\mathrm T}H_{\mathrm T}$}
\end{phflow}

\phboxhead{物理量・防護量・実用量}
ここで、3つの量のちがいを整理しておきましょう（環境省の基礎資料の分け方）。
\begin{bkfigure}
\begin{tikzpicture}[x=1mm, y=1mm]
  \node[txo box, fill=bkBlue!8, minimum width=36mm, minimum height=15mm] (b) at (18,30) {\textbf{物理量}\\放射能 Bq\\吸収線量 Gy};
  \node[txo box, fill=bkOrange!10, minimum width=36mm, minimum height=15mm] (h) at (76,30) {\textbf{防護量}\\等価線量 Sv\\実効線量 Sv};
  \node[txo box, fill=bkGreen!10, minimum width=40mm, minimum height=15mm] (j) at (76,4) {\textbf{実用量}\\1 cm 線量当量 Sv\\（測定器が示す値）};
  \draw[txo arr, bkInk] (b) -- node[txo lab, above, align=center] {$w_{\mathrm R}$・$w_{\mathrm T}$\\で重み} (h);
  \draw[txo arr, bkInk] (b.south) |- node[txo lab, below, pos=0.75] {はかれる量から定義} (j.west);
  \draw[txo arr, bkInk, dashed] (j) -- node[txo lab, right, xshift=1mm, align=left] {安全側に\\少し大きく近似} (h);
  \node[txo lab, anchor=west, align=left] at (96,30) {人の臓器の線量から\\計算する量。直接は\\はかれない};
\end{tikzpicture}
\bkfigcaption{物理量・防護量・実用量。実効線量（防護量）は直接はかれないので、測定器は、はかれる物理量から定義した実用量（1 cm 線量当量など）を Sv で表示する。実用量は、実効線量より少し大きな値になるように決められている。}\label{ph:原-14:fig-3ryo}
\end{bkfigure}
\textbf{物理量}（Bq・Gy）は、測定器や計算で直接決められる量です。\textbf{防護量}（等価線量・実効線量）は、人を守るための管理に使う量で、体の中の臓器ごとの線量から計算するため、直接ははかれません。そこで、空間の線量率をはかるサーベイメーターなどは、体の表面から 1 cm の深さの線量などで定義した\textbf{実用量}を、Sv の単位で表示しています。実用量は、防護量を安全側に（少し大きめに）見積もるように作られています。ニュースの「\si{\micro Sv/h}」は、ふつうこの実用量です。

\Needspace{8\baselineskip}\bksubsection{6}{点線源の逆2乗の法則と、検出器の立体角}
\label{ph:原-14:sec6}
\noindent\begin{minipage}[c]{0.55\linewidth}\raggedright
大きさが無視できるほど小さな放射線の出どころ（\textbf{点線源}）から、1秒に $S$ 個の光子が、四方八方に同じように出ているとします。途中で何にもぶつからなければ、光子は距離 $r$ の球面（面積 $4\pi r^2$）を全部通りぬけるので、1秒に 1 cm$^2$ を通る光子の数は
\[ \frac{S}{4\pi r^2} \]
です。線量率（1時間あたりの線量）は、この数に比例するので、\textbf{距離の2乗に反比例}します。音の逆2乗の法則（波-7 ②）と同じです。
\end{minipage}\hfill
\begin{minipage}[c]{0.42\linewidth}\centering
\begin{tikzpicture}[x=1mm, y=1mm]
  \fill[bkRed!80!black] (0,0) circle (1.2);
  \draw[bkGray, line width=0.6pt] (0,0) circle (12);
  \draw[bkGray, line width=0.6pt] (0,0) circle (24);
  \foreach \a in {0,30,...,330} \draw[bkOrange!85!black, line width=0.5pt, -{Stealth[length=1.4mm]}] (\a:1.8) -- (\a:28);
  \fill[bkBlue!40] (12,-2) rectangle (13,2);
  \fill[bkBlue!40] (24,-4) rectangle (25,4);
  \node[txo lab, fill=white, inner sep=1pt] at (-10,-17) {$r$};
  \node[txo lab, fill=white, inner sep=1pt] at (-15,-25) {$2r$};
  \node[txo lab, anchor=north, align=center] at (0,-29) {同じ数の線が、2倍の距離では\\4倍の面積に広がる};
\end{tikzpicture}
\end{minipage}

\phboxhead{線量率定数}
γ線を出す核種ごとに、「放射能 1 MBq の点線源から 1 m の所の線量率」を表にしたものを、\textbf{線量率定数} $\Gamma$（ガンマ）といいます。\bkidx{せんりょうりつていすう}{線量率定数}放射能 $A$〔MBq〕の点線源から距離 $r$〔m〕の所の線量率は、
\pheq{\dot H=\frac{\Gamma A}{r^2}}{\si{\micro Sv/h} ＝ $\Gamma$ × MBq ÷ m$^2$}
\label{ph:原-14:gyaku2}です（$\dot H$ の点は「1時間あたり」）。セシウム137 の 1 cm 線量当量率定数は
\[ \Gamma=\qty{0.0927}{\micro Sv.m^2/(MBq.h)} \]
です（日本アイソトープ協会『アイソトープ手帳』の値）。1 MBq の点線源から 1 m の所で約 \qty{0.093}{\micro Sv/h}、50 cm では4倍の約 \qty{0.37}{\micro Sv/h}、2 m では $\tfrac14$ の約 \qty{0.023}{\micro Sv/h} です（例題5）。

\phboxhead{線量率定数を、物理の量から見積もる}
$\Gamma$ の大きさは、原-13 で学んだことから見積もれます。セシウム137 は、こわれる100回のうち約85回、$E_\gamma=\qty{0.662}{MeV}$ のγ線を出します。1 MBq（1秒に $10^6$ 回こわれる）の点線源から 1 m（100 cm）の所を、1秒に 1 cm$^2$ あたり通る光子の数は、
\[ \frac{10^6\times0.85}{4\pi\times100^2}=6.8\ \text{個} \]
です。空気 1 g が、この光子から受け取るエネルギーの割合は、空気の質量エネルギー吸収係数（NIST の表で約 \qty{0.029}{cm^2/g}）で決まります。計算すると、空気が1時間に受け取るエネルギーは、空気 1 kg あたり約 $\qty{0.076}{\micro J}$、つまり約 \qty{0.076}{\micro Gy/h} になります（物理量）。$\Gamma=0.093$ はこれより約2割大きい値ですが、これは、実用量の 1 cm 線量当量が、体の中で散乱した光子の分などをふくめて、安全側に大きめに決められているからです。

\phboxhead{検出器が見る割合：立体角}
\noindent\begin{minipage}[c]{0.56\linewidth}\raggedright
半径 $a$ の円い窓の検出器を、点線源から距離 $d$ の所に、窓を線源に向けて置きます。線源から出る放射線のうち、窓に飛びこむ割合は、窓が線源から見て占める立体角 $\Omega$（8-8⑬）の、まわり全体 $4\pi$ に対する割合です。窓は半分の角 $\alpha$（$\cos\alpha=d/\sqrt{d^2+a^2}$）の円すいを切りとるので、
\end{minipage}\hfill
\begin{minipage}[c]{0.41\linewidth}\centering
\begin{tikzpicture}[x=1mm, y=1mm]
  \fill[bkRed!80!black] (0,0) circle (1.1);
  \fill[bkYellow!35] (0,0) -- (36,9) -- (36,-9) -- cycle;
  \draw[bkInk, line width=0.6pt] (0,0) -- (36,9) (0,0) -- (36,-9);
  \filldraw[fill=bkBlue!40, draw=bkInk, line width=0.6pt] (36,-9) rectangle (40,9);
  \draw[bkGray, line width=0.5pt, dash pattern=on 1.5pt off 1pt] (0,0) -- (36,0);
  \draw[bkInk, line width=0.5pt] (10,0) arc (0:14:10);
  \node[txo lab] at (13,1.7) {$\alpha$};
  \draw[txo dim] (0,-12) -- (36,-12);
  \node[txo lab, anchor=north] at (18,-12.5) {$d$};
  \draw[txo dim] (43,0) -- (43,9);
  \node[txo lab, anchor=west] at (43.5,4.5) {$a$};
  \node[txo lab, anchor=south] at (0,2) {線源};
  \node[txo lab, anchor=south] at (38,9.5) {検出器};
\end{tikzpicture}
\end{minipage}
\pheq{\frac{\Omega}{4\pi}=\frac{2\pi(1-\cos\alpha)}{4\pi}=\frac12\Bigl(1-\frac{d}{\sqrt{d^2+a^2}}\Bigr)}{1（割合）}
\label{ph:原-14:rittai}です。$d$ が $a$ よりずっと大きいときは、窓の面積 $\pi a^2$ を球面の面積 $4\pi d^2$ でわった $\dfrac{a^2}{4d^2}$（逆2乗の法則）に近づきます。半径 2 cm の窓を 10 cm の所に置くと、正確には 0.97 \%、逆2乗の近似では 1.0 \% で、ほぼ同じですが、5 cm まで近づけると 3.6 \% と 4.0 \% で、1割ずれます（例題11）。

\begin{phmodel}[点線源]
\phignore{線源の大きさ。空気による減弱（0.662 MeV のγ線は、10 m 進むと約1割減る）。床や壁で散乱して届く放射線。}
\phassume{線源は点で、放射線はどの向きにも同じように出て、まっすぐ進む。}
\phlimit{線源の大きさと同じくらいまで近づいたとき（例題11）。地面に広く散らばった放射性物質（このときは距離を2倍にしても $\tfrac14$ にはならない）。}
\end{phmodel}

\Needspace{8\baselineskip}\bksubsection{7}{外からの被ばくを減らす3原則：時間・距離・遮へい}
\label{ph:原-14:sec7}
体の外にある線源から受ける線量（\textbf{外部被ばく}）は、線量率 × 時間で決まります。⑥と原-13 ③を組み合わせると、点線源から距離 $r$ の所に、厚さ $x$ の遮へいごしに、時間 $t$ だけいるときの線量は
\pheq{H=\frac{\Gamma A}{r^2}\,t\,e^{-\mu x}}{\si{\micro Sv} ＝ \si{\micro Sv/h} × h}
\label{ph:原-14:3gensoku}です（ビルドアップは無視：原-13 ⑩）。この式から、線量を減らす方法は3つあることがわかります。これを\textbf{外部被ばくを減らす3原則}といいます。\bkidx{がいぶひばくをへらす3げんそく}{外部被ばくを減らす3原則}
\begin{bkfigure}
\begin{tikzpicture}[x=1mm, y=1mm]
  \foreach \i/\t/\s in {0/時間/{いる時間を短く\\線量は時間に比例}, 44/距離/{線源から離れる\\距離2倍で $\tfrac14$}, 88/遮へい/{間に物を置く\\半価層ごとに $\tfrac12$}} {
    \filldraw[fill=bkTheme!6, draw=bkTheme!70, line width=0.7pt, rounded corners=2mm] (\i,0) rectangle (\i+40,30);
    \node[font=\sffamily\bfseries\small, text=bkDeep] at (\i+20,25) {\t};
    \node[txo lab, align=center] at (\i+20,6) {\s};
  }
  % 時計
  \draw[bkInk, line width=0.7pt] (20,15.5) circle (4.5);
  \draw[bkInk, line width=0.8pt] (20,15.5) -- (20,19) (20,15.5) -- (22.5,15.5);
  % 距離
  \fill[bkRed!80!black] (52,15.5) circle (1.1);
  \draw[txo dim] (54,15.5) -- (74,15.5);
  \node[font=\sffamily\scriptsize] at (78,15.5) {人};
  % 遮へい
  \fill[bkRed!80!black] (95,15.5) circle (1.1);
  \fill[bkGray!55] (105,12.5) rectangle (110,21);
  \draw[bkOrange!85!black, line width=0.8pt, -{Stealth[length=1.5mm]}] (97,15.5) -- (104.5,15.5);
  \draw[bkOrange!85!black, line width=0.5pt, dash pattern=on 1.5pt off 1pt, -{Stealth[length=1.5mm]}] (110.5,15.5) -- (118,15.5);
\end{tikzpicture}
\bkfigcaption{外部被ばくを減らす3原則（環境省の基礎資料 4.3節）。}\label{ph:原-14:fig-3gensoku}
\end{bkfigure}
\begin{itemize}[label={\color{bkTheme}●}, leftmargin=1.6\zw, itemsep=1pt, topsep=1pt]
\item \textbf{時間}：線量は時間に比例します。時間を半分にすれば、線量も半分です。
\item \textbf{距離}：線量は距離の2乗に反比例します。距離を2倍にすれば $\tfrac14$、3倍にすれば $\tfrac19$ です。
\item \textbf{遮へい}：半価層の厚さの物を1枚置くごとに $\tfrac12$ です。建物の中にいることも遮へいになります。地面にたまった放射性物質からの線量は、木造の家の中で外の約4割になるとされています（環境省の基礎資料）。
\end{itemize}
3つのうち、どれか1つだけでなく、組み合わせることで、大きく減らせます（例題6）。

\begin{phbaisuru}{線源からの距離 $r$}
\noindent\begin{minipage}[c]{0.46\linewidth}\raggedright
線量率は $\tfrac14$ になります。3倍で $\tfrac19$、4倍で $\tfrac1{16}$ です。線源の近くでは、少し離れるだけで大きく減るので、線源を手で直接持たず、長いトングなどを使って作業します。グラフは、セシウム137 の 1 MBq の点線源（$\Gamma A=\qty{0.093}{\micro Sv.m^2/h}$）の線量率です。
\end{minipage}\hfill
\begin{minipage}[c]{0.50\linewidth}\centering
\begin{tikzpicture}
\begin{axis}[bkplot, axis lines=left, width=50mm, height=40mm, xmin=0, xmax=3.3, ymin=0, ymax=0.42, xtick={0.5,1,2,3}, ytick={0.1,0.2,0.3,0.4}, xlabel={$r$（m）}, ylabel={\si{\micro Sv/h}}]
  \addplot[bkcurve, domain=0.47:3.2, samples=80] {0.0927/x^2};
  \fill[bkBlue] (axis cs:0.5,0.371) circle (1.5pt) (axis cs:1,0.0927) circle (1.5pt) (axis cs:2,0.0232) circle (1.5pt);
  \node[font=\sffamily\footnotesize, anchor=west, inner sep=1pt] at (axis cs:0.6,0.37) {0.37};
  \node[font=\sffamily\footnotesize, anchor=south west, inner sep=1pt] at (axis cs:1.05,0.10) {0.093};
  \node[font=\sffamily\footnotesize, anchor=south west, inner sep=1pt] at (axis cs:2.05,0.03) {0.023};
\end{axis}
\end{tikzpicture}
\end{minipage}
\end{phbaisuru}

\Needspace{8\baselineskip}\bksubsection{8}{体の中からの被ばく：内部被ばくと食品の基準}
\label{ph:原-14:sec8}
放射性物質を、食べ物や飲み物、空気といっしょに体の中に取りこむと、体の中から放射線を受けます。これを\textbf{内部被ばく}といいます。\bkidx{ないぶひばく}{内部被ばく}体の外からは皮ふで止まる α線や β線も、体の中では、すぐそばの細胞にエネルギーを渡します（原-13 ⑤）。

取りこんだ物は、時間とともに、こわれて（物理学的半減期）、また尿や便といっしょに出ていって（生物学的半減期）、減っていきます。そこで、1回取りこんだ量から、将来にわたって受ける線量の合計（大人は50年、子どもは70歳まで）を計算し、それを取りこんだ年に受けたものとみなします。これを\textbf{預託実効線量}といいます。\bkidx{よたくじっこうせんりょう}{預託実効線量}1 Bq を取りこんだときの預託実効線量を\textbf{線量係数}（実効線量係数）とよび、核種・取りこみ方・年齢ごとに決められています。
\pheq{E=(\text{取りこんだ放射能〔Bq〕})\times(\text{線量係数〔\si{\micro Sv}/Bq〕})}{\si{\micro Sv} ＝ Bq × \si{\micro Sv}/Bq}
\label{ph:原-14:naibu}
\begin{bkcenter}
\begin{bktable}{colspec={Q[l,m,wd=25mm] Q[c,m,wd=16mm] Q[c,m,wd=15mm] Q[c,m,wd=16mm] Q[c,m,wd=16mm] Q[c,m,wd=18mm]}, rowsep=1pt, row{1}={bg=bkTheme!10}, cells={font=\small}}
核種 & ストロンチウム90 & ヨウ素131 & セシウム134 & セシウム137 & トリチウム（水） \\
係数〔\si{\micro Sv}/Bq〕 & 0.028 & 0.022 & 0.019 & 0.013 & 0.000018 \\
\end{bktable}
\bktabcaption{口から取りこんだときの線量係数（大人。ICRP Publication 119 による。環境省の基礎資料）。子どもは別の値が決められている。}\label{ph:原-14:tab-keisu}
\end{bkcenter}
トリチウム（三重水素）は、β線のエネルギーがとても小さく、水として体に入っても、約10日で半分が体の外に出ていくので（生物学的半減期）、線量係数がセシウム137 の約700分の1と小さくなっています。

日本では、食品にふくまれる放射性セシウムの\textbf{基準値}を、一般食品 1 kg あたり 100 Bq、牛乳と乳児用食品 50 Bq、飲料水 10 Bq と決めています（2012年4月から）。\bkidx{きじゅんち}{基準値}どんな食品を食べても、食品からの線量が1年に 1 mSv をこえないように、余裕をもって決めた値です。たとえば、100 Bq/kg のセシウム137 をふくむ食品を 1 kg 食べたときの預託実効線量は $100\times0.013=\qty{1.3}{\micro Sv}$ です（例題8）。

\Needspace{8\baselineskip}\bksubsection{9}{「ある」と言える最小の量：検出限界}
\label{ph:原-14:sec9}
測定器は、自然放射線（⑩）の分を、いつも少し数えています。これを\textbf{バックグラウンド}といいます。食品などの試料をはかって、数がバックグラウンドより少し多かったとしても、それが試料の放射線なのか、バックグラウンドのばらつき（③）なのかは、すぐには決められません。

そこで、バックグラウンドの平均に、そのばらつきの標準偏差 $\sigma$ の3倍を足した値をこえたときに、「試料に放射性物質がある（検出された）」と判断します。これを\textbf{3σ 法}といい、\bkidx{3しぐまほう}{3σ法}その境目の量を\textbf{検出限界}といいます。\bkidx{けんしゅつげんかい}{検出限界}バックグラウンドだけで、ばらつきのせいで 3σ をこえてしまう確率は、約 0.1 \% しかありません（3-17 の正規分布）。

測定の結果が「不検出（ND）」とあるのは、「放射性物質が0」という意味ではなく、「検出限界より少ない」という意味です。検出限界は、③のとおり、\textbf{長くはかるほど、試料が多いほど、小さく}なります。背景が 30 cpm（1分に30カウント）の測定器なら、1分だけはかると $3\sigma$ は約 16 cpm ですが、10分はかると約 5.2 cpm、40分はかると約 2.6 cpm になります（例題10）。

\Needspace{8\baselineskip}\bksubsection{10}{自然放射線と医療の線量：量で比べる}
\label{ph:原-14:sec10}
私たちは、宇宙から来る放射線（宇宙線）、大地にふくまれるウランやトリウムなどからの放射線、空気中のラドン、食べ物にふくまれるカリウム40 や鉛210・ポロニウム210 などから、いつも放射線を受けています。これを\textbf{自然放射線}といいます。\bkidx{しぜんほうしゃせん}{自然放射線}日本人が1年に自然から受ける実効線量は、平均で約 \qty{2.1}{mSv}（世界の平均は約 \qty{2.4}{mSv}）です。
\begin{bkfigure}
\begin{tikzpicture}
\begin{axis}[bkplot, axis lines=left, width=72mm, height=54mm, xmin=0, xmax=0.95, ymin=0.4, ymax=6.6, ytick={1,2,3,4,5,6}, yticklabels={その他（炭素14・温泉・飛行機など）,大地（外から）,宇宙線（外から）,ラドンなど（吸う）,カリウム40（食べる）,鉛210・ポロニウム210（食べる）}, xtick={0,0.2,0.4,0.6,0.8}, xlabel={mSv/年}, grid=none, every axis x label/.style={at={(axis description cs:0.5,-0.18)}, anchor=north, font=\sffamily\small}, y tick label style={font=\sffamily\footnotesize}]
  \fill[bkBlue!60] (axis cs:0,0.7) rectangle (axis cs:0.027,1.3);
  \fill[bkBlue!60] (axis cs:0,1.7) rectangle (axis cs:0.33,2.3);
  \fill[bkBlue!60] (axis cs:0,2.7) rectangle (axis cs:0.30,3.3);
  \fill[bkBlue!60] (axis cs:0,3.7) rectangle (axis cs:0.47,4.3);
  \fill[bkBlue!60] (axis cs:0,4.7) rectangle (axis cs:0.18,5.3);
  \fill[bkBlue!60] (axis cs:0,5.7) rectangle (axis cs:0.80,6.3);
  \node[font=\sffamily\footnotesize, anchor=west] at (axis cs:0.05,1) {0.03};
  \node[font=\sffamily\footnotesize, anchor=west] at (axis cs:0.35,2) {0.33};
  \node[font=\sffamily\footnotesize, anchor=west] at (axis cs:0.32,3) {0.30};
  \node[font=\sffamily\footnotesize, anchor=west] at (axis cs:0.49,4) {0.47};
  \node[font=\sffamily\footnotesize, anchor=west] at (axis cs:0.20,5) {0.18};
  \node[font=\sffamily\footnotesize, anchor=west] at (axis cs:0.82,6) {0.80};
\end{axis}
\end{tikzpicture}
\bkfigcaption{日本人が自然から1年に受ける実効線量の内訳（合計 約 2.1 mSv。環境省の基礎資料 2.5節、原子力安全研究協会「生活環境放射線 第3版 増補版」（2024年）による）。魚や貝をよく食べるため、鉛210・ポロニウム210 が多いのが日本の特徴。}\label{ph:原-14:fig-shizen}
\end{bkfigure}
これを1時間あたりにすると、$2.1\times10^3\div(365.25\times24)\approx\qty{0.24}{\micro Sv/h}$ です（例題7）。このうち、体の外から受ける分（宇宙線と大地：約 0.63 mSv/年）は約 \qty{0.07}{\micro Sv/h} で、ふつうの場所の空間線量率（環境省の資料では、生活空間のほとんどで1時間に 0.01〜1 \si{\micro Sv}、測定器のバックグラウンドは 0.1 \si{\micro Sv/h} 程度）と、桁が合っています。

日本では、CT 検査などの医療で受ける線量が、1人あたり平均で年に約 \qty{2.6}{mSv} と見積もられていて、自然放射線と同じくらいの大きさです。医療の放射線は、検査や治療で得られる利益が、放射線による害より大きいときに使われます。

\Needspace{8\baselineskip}\bksubsection{11}{体への影響：確定的影響と確率的影響}
\label{ph:原-14:sec11}
放射線による体への影響は、2種類に分けて考えます。
\begin{itemize}[label={\color{bkTheme}●}, leftmargin=1.6\zw, itemsep=1pt, topsep=1pt]
\item \textbf{確定的影響}（組織反応）：たくさんの細胞が死ぬことで起きる影響（皮ふが赤くなる、血をつくるはたらきが落ちる、など）。ある線量（\textbf{しきい線量}）より少なければ起きず、それをこえると、線量が多いほど症状が重くなる。\bkidx{かくていてきえいきょう}{確定的影響}
\item \textbf{確率的影響}：DNA の傷が残った細胞から、何年もたってがんなどが起きる影響。線量が多いほど起きる\textbf{確率}が高くなるが、症状の重さは線量によらない。\bkidx{かくりつてきえいきょう}{確率的影響}
\end{itemize}

\begin{bkcenter}
\begin{bktable}{colspec={Q[l,m,wd=30mm] Q[l,m,wd=34mm] Q[c,m,wd=34mm]}, rowsep=1pt, row{1}={bg=bkTheme!10}, cells={font=\small}}
影響 & 臓器・組織 & しきい値〔Gy〕 \\
一時的な不妊 & 精巣 & 約 0.1 \\
血をつくるはたらきの低下 & 骨髄 & 約 0.5 \\
白内障（視力の低下） & 眼 & 約 0.5 \\
一時的な脱毛 & 皮ふ & 約 4 \\
皮ふが赤くなる & 皮ふ（広い範囲） & 3〜6 以下 \\
\end{bktable}
\bktabcaption{γ線を短い時間に受けたときの、確定的影響のしきい値の例（1 \% の人に症状が出る線量。ICRP 2007年勧告などによる。環境省の基礎資料 3.3節）。}\label{ph:原-14:tab-shikii}
\end{bkcenter}

\begin{bkfigure}
\begin{tikzpicture}
\begin{axis}[bkplot, axis lines=left, name=ka, width=50mm, height=38mm, xmin=0, xmax=6, ymin=0, ymax=1.15, xtick={1.1}, xticklabels={しきい線量}, ytick=\empty, xlabel={線量}, ylabel={起きる割合}]
  \addplot[bkcurve, domain=0:5.8, samples=80] {1/(1+exp(-3*(x-2.6)))};
  \draw[bkGray, dash pattern=on 1.5pt off 1pt, line width=0.5pt] (axis cs:1.1,0) -- (axis cs:1.1,0.5);
  \node[font=\sffamily\footnotesize, anchor=north, inner sep=1pt] at (axis cs:3,-0.25) {確定的影響};
\end{axis}
\begin{axis}[bkplot, axis lines=left, at={(ka.east)}, anchor=west, xshift=20mm, width=50mm, height=38mm, xmin=0, xmax=6, ymin=0, ymax=1.15, xtick={2}, xticklabels={100 mSv}, ytick=\empty, xlabel={線量}, ylabel={確率の増え方}]
  \addplot[bkcurve, domain=2:5.8] {0.15*x};
  \addplot[bkcurve, dashed, domain=0:2] {0.15*x};
  \addplot[bkcurve2, densely dotted, domain=0:1] {0};
  \addplot[bkcurve2, densely dotted, domain=1:2] {0.3*(x-1)};
  \node[font=\sffamily\footnotesize, anchor=north, inner sep=1pt] at (axis cs:3,-0.25) {確率的影響};
  \node[font=\sffamily\footnotesize, anchor=north west, inner sep=1pt] at (axis cs:0.1,0.95) {LNT（破線）};
\end{axis}
\end{tikzpicture}
\bkfigcaption{確定的影響（左）と確率的影響（右）の線量との関係（模式図）。右の 100 mSv より上（実線）では、線量に比例してがんの確率が増えることが、原爆の被爆者の調査などからわかっている。それより下は見分けがむずかしく、そのまま直線をのばす見方（LNT、破線）と、ある線量より下では増えないとする見方（点線の例）がある。}\label{ph:原-14:fig-eikyo}
\end{bkfigure}

\phboxhead{100 mSv あたり 0.5 \%}
国際放射線防護委員会（ICRP）は、子どもから大人までをふくむ集団では、受けた線量 100 mSv あたり、がんで亡くなる確率が 0.5 \% 増えるとして、防護を考えることにしています。これは、広島・長崎の原爆の被爆者の長い調査のデータをもとに、低い線量率での値を見積もったものです。日本では、もともと約30 \% の人ががんで亡くなっています。1000人の集団全員が 100 mSv を受けたとすると、がんで亡くなる人は、計算の上では300人から305人に増えることになります（例題13）。けれども、もとの300人も、年や地域で変わるので、この5人の差を実際に見分けるのは、とてもむずかしいと考えられています。

\begin{phmada}[100 mSv より少ない線量で、がんは増えるか]
100 mSv より少ない線量でのがんのリスクについては、\textbf{科学的な決着がついていません}。米国科学アカデミーの委員会（2006年）や米国放射線防護審議会（2018年）は、低い線量でも線量に比例してリスクがあるとする\textbf{直線しきい値なしモデル}（LNT モデル）を支持する見解を出しています。一方、フランスの医学アカデミーと科学アカデミー（2005年）は、ある線量より低ければがんは実際には起きず、LNT モデルは大きすぎる見積もりだ、という見解を出しています。自然放射線が多い地域の住民の調査や、細胞が傷を直すしくみの研究などが、その根拠です。低い線量では、放射線による増加が、ほかの原因によるがんのばらつきよりずっと小さいので、疫学の調査で見分けることがむずかしいのです。

\textbf{放射線防護では、LNT モデルを使います}。それが正しいと決まったからではなく、わからない所を安全側に見積もり、管理を単純で合理的にするためです。ICRP は同時に、「ごく小さい線量を多くの人が受けたとして、がんの数を計算するのは、公衆の健康の計画には適切でない」とも述べています。
\end{phmada}

\Needspace{8\baselineskip}\bksubsection{12}{線量限度と、放射線を使う仕事}
\label{ph:原-14:sec12}
放射線を仕事で使う人（放射線業務従事者）が受ける実効線量は、法令（電離放射線障害防止規則 第4条）で、\textbf{5年間で 100 mSv をこえず、かつ1年間で 50 mSv をこえない}ように決められています。妊娠する可能性のある女性は、3か月で 5 mSv です。一般の人（公衆）については、ICRP は、管理された線源から受ける線量を1年に 1 mSv と勧告しています（医療と自然放射線は数えない）。\bkidx{せんりょうげんど}{線量限度}

\begin{bkcenter}
\begin{bktable}{colspec={Q[l,m,wd=56mm] Q[l,m,wd=54mm]}, rowsep=1pt, row{1}={bg=bkTheme!10}, cells={font=\small}}
決まり & 値 \\
放射線業務従事者の実効線量 & 5年で 100 mSv かつ 1年で 50 mSv \\
女性の放射線業務従事者（妊娠の可能性がある人） & 3か月で 5 mSv \\
管理区域（標識で示す区域） & 3か月で 1.3 mSv をこえるおそれがある所 \\
一般公衆（ICRP の勧告） & 1年で 1 mSv（医療・自然は除く） \\
\end{bktable}
\bktabcaption{線量限度など（電離放射線障害防止規則 第3条・第4条、ICRP 2007年勧告。環境省の基礎資料 4.1節・4.2節）。}\label{ph:原-14:tab-gendo}
\end{bkcenter}

ここで大切なのは、\textbf{線量限度は、安全と危険の境界線ではない}ということです。線量限度は、管理する線源からの線量の合計が、その値をこえないように管理するための基準です。こえなければいくら受けてもよい、という値ではなく、限度より下でも、できるだけ線量を下げる努力（\textbf{防護の最適化}）が求められます。また、1 mSv を少しこえたら急に危険になる、という値でもありません。

\phboxhead{放射線を使う仕事：放射線透過試験}
工場や建設の現場では、X線やγ線で溶接の内部を調べる\textbf{放射線透過試験}（原-13 ⑪）が行われます。作業する人は、個人線量計をつけて自分の線量を記録し、線源のまわりに立ち入りを禁止する区域をつくり、3原則（⑦）で線量を減らします。作業の計画では、例題12 のように、「その場所の線量率で、何時間までなら作業できるか」を計算して決めます。

\begin{phsaki}{放射線計測と放射線防護の先へ}
この章の測定器は、放射線の「数」と「エネルギー」をはかりました。専門の本では、検出器の中で放射線がエネルギーを落とすようすを計算機でまねて、測定器の感度を決めます（原-13 のモンテカルロ法）。体の中の臓器ごとの線量は、人の体の形をこまかく再現した計算の模型で求めます。体の中にある放射性物質は、体の外から γ線をはかる装置（ホールボディカウンター）で調べます。これらは、放射線計測学・保健物理学（放射線防護学）で学びます。
\end{phsaki}

\begin{senmonsho}[放射線防護の専門書の書き方]
\begin{bkhonbun}
外部被ばくの管理には実用量として周辺線量当量 $H^*(10)$ および個人線量当量 $H_{\mathrm p}(10)$ を用いる．なお旧単位との換算は $1\ \mathrm{Ci}=3.7\times10^{10}\ \mathrm{Bq}$，$1\ \mathrm{rad}=10^{-2}\ \mathrm{Gy}$，$1\ \mathrm{rem}=10^{-2}\ \mathrm{Sv}$，照射線量 $1\ \mathrm R=2.58\times10^{-4}\ \mathrm{C/kg}$ である．
\end{bkhonbun}
\bkyaku{$H^*(10)$ は「エイチ・スター・テン」と読み、この章の「1 cm 線量当量」（深さ 10 mm）のこと。$H_{\mathrm p}(10)$ は個人線量計で使う実用量。Ci（キュリー）・rad（ラド）・rem（レム）・R（レントゲン）は、古い本や外国の資料に出てくる昔の単位。この本では使わない。}
\end{senmonsho}

\begin{taisetsu}
\def\tsl#1{\makebox[7.6\zw][l]{\sffamily\bfseries\small\color{\bkTT}#1}}
\tsl{数える}$\sigma\approx\sqrt N$、数え落とし $n=\dfrac{m}{1-m\tau}$、検出限界は背景 $+3\sigma$\\[3pt]
\tsl{線量}$D=\dfrac{\Delta E_{\text{吸}}}{\Delta m}$〔Gy〕、$H_{\mathrm T}=w_{\mathrm R}D_{\mathrm T}$、$E=\textstyle\sum w_{\mathrm T}H_{\mathrm T}$〔Sv〕\\[3pt]
\tsl{点線源}$\dot H=\dfrac{\Gamma A}{r^2}$、$\dfrac{\Omega}{4\pi}=\dfrac12\Bigl(1-\dfrac{d}{\sqrt{d^2+a^2}}\Bigr)$\\[3pt]
\tsl{3原則}$H=\dfrac{\Gamma A}{r^2}\,t\,e^{-\mu x}$（時間・距離・遮へい）\\[3pt]
\tsl{内部被ばく}預託実効線量＝取りこんだ Bq × 線量係数\\[3pt]
\tsl{量で比べる}自然 約 2.1 mSv/年（日本）、医療 約 2.6 mSv/年、100 mSv あたり 0.5 \%
\end{taisetsu}

\begin{kotoba}
\small
\bkword{吸収線量}{物 1 kg が放射線から受け取ったエネルギー。単位 Gy（グレイ）＝J/kg。}
\bkword{等価線量}{吸収線量に、放射線の種類による重み $w_{\mathrm R}$ をかけた量。単位 Sv。}
\bkword{実効線量}{等価線量に、組織ごとの重み $w_{\mathrm T}$ をかけて全身で合計した量。単位 Sv。}
\bkword{実用量}{測定器が表示する、実効線量を安全側に見積もるための量（1 cm 線量当量など）。}
\bkword{不感時間}{検出器が1つ数えたあと、次を数えられない短い時間。}
\bkword{検出限界}{「ある」と判断できる最小の量。背景の平均 ＋ 3σ（3σ 法）。}
\bkword{預託実効線量}{体に取りこんだ放射性物質から、将来にわたって受ける線量の合計。}
\bkword{確定的影響}{しきい線量をこえると起きる影響。線量が多いほど重くなる。}
\bkword{確率的影響}{がんなど、線量が多いほど起きる確率が高くなる影響。}
\bkword{線量限度}{管理する線源からの線量が、こえてはいけない値。安全と危険の境界線ではない。}
\end{kotoba}

\bkkigocards{\bkkigocard{$D$}{ディー}{吸収線量〔Gy〕。$D=\Delta E_{\text{吸}}/\Delta m$。原-9 の拡散係数 $D_{\mathrm n}$ とは別の量}{absorbed dose}{原-14}}
{\bkkigocard{$H_{\mathrm T}$}{エイチ・ティー}{等価線量〔Sv〕。$H_{\mathrm T}=w_{\mathrm R}D_{\mathrm T}$。T は組織（tissue）}{equivalent dose}{原-14}}
\bkkigocards{\bkkigocard{$E$}{イー}{実効線量〔Sv〕。$E=\sum w_{\mathrm T}H_{\mathrm T}$。原-14 では、エネルギーは $E_\gamma$ などと添字で書く}{effective dose}{原-14}}
{\bkkigocard{$w_{\mathrm R}$}{ダブリュー・アール}{放射線加重係数。放射線の種類による単位のない重み（α線 20 など）}{radiation weighting factor}{原-14}}
\bkkigocards{\bkkigocard{$w_{\mathrm T}$}{ダブリュー・ティー}{組織加重係数。組織ごとの単位のない重みで、合計は1}{tissue weighting factor}{原-14}}
{\bkkigocard{$\tau$}{タウ}{不感時間〔s〕。数え落とし $n=m/(1-m\tau)$。6-4・原-12 の時定数 $\tau$ とは別の量}{dead time}{原-14}}
\bkkigocards{\bkkigocard{$\Gamma$}{ガンマ（大文字）}{線量率定数〔\si{\micro Sv}·m$^2$/(MBq·h)〕。1 MBq の点線源から 1 m の線量率}{dose rate constant}{原-14}}{}

%% ======================== 例題 ========================
\Needspace{10\baselineskip}\bksection{reidai}{例題でたしかめよう}
\label{ph:原-14:ex1}
\begin{phreidai}[GM 計数管の数え落とし（はてな？）]
GM 計数管で1分はかったら、\num{30000} カウントでした。不感時間を \qty{100}{\micro s} とすると、本当に来た放射線は1秒に何個ですか。数え落としは何 \% ですか。
\phbamen
強い線源の近くで、GM 計数管の表示を正しく読みとる場面です。
\phmushi
不感時間は、いつも同じ長さで、そのあいだに来た放射線は、数えられないだけで、時間をのばさないとします。
\phzu
数えた数 $m=\num{30000}\div60=500$ /s、$\tau=\qty{1.0e-4}{s}$。
\phriyu
1秒のうち $m\tau$ 秒が「休み時間」なので、数え落としの式（③）を使います。
\phtateru
\[ n=\frac{m}{1-m\tau} \]
\phtoku
\[ m\tau=500\times1.0\times10^{-4}=0.050,\qquad n=\frac{500}{1-0.050}=\frac{500}{0.95}=526\ /\text{s} \]
数え落としは $(526-500)\div526=0.050$ で、5 \% です。
\phtashikame
単位：(1/s) ÷ 1＝1/s。\\
大きさ：数え落としの割合は $m\tau$ に等しく、5 \%。\\
極端な場合：$m$ が小さい（弱い線源）なら、$n\approx m$ で補正はいりません。
\kotae 1秒に約 526 個（数え落とし 5 \%）
\end{phreidai}
\begin{jibun}
同じ GM 計数管で、1秒に 2000 カウントのとき、本当の数と数え落としの割合を求めましょう。
\bkjans{$n=2000/(1-0.20)=\num{2500}$ /s、数え落とし 20 \%}
\end{jibun}

\label{ph:原-14:ex2}
\begin{phreidai}[1 Gy は、どれくらいのエネルギーか]
体重 60 kg の人の全身が、1 Gy の吸収線量を受けました。体が受け取ったエネルギーと、それが全部熱になったときの体温の上がり方を求めましょう。体の比熱は水と同じ $\qty{4.18e3}{J/(kg.K)}$ とします。
\phbamen
放射線の「エネルギーの量」を、身近な熱の量と比べる場面です。
\phmushi
エネルギーは全身に均一に吸収されるとし、熱が体の外へ逃げることは無視します。
\phzu
$D=\qty{1}{Gy}$、$m=\qty{60}{kg}$、比熱 $c=\qty{4.18e3}{J/(kg.K)}$。
\phriyu
吸収線量の定義（④）からエネルギーを求め、熱のバランス（熱-1 ③）で温度の上がり方を求めます。
\phtateru
\[ \Delta E_{\text{吸}}=Dm,\qquad \Delta T=\frac{\Delta E_{\text{吸}}}{cm} \]
\phtoku
\[ \Delta E_{\text{吸}}=1\times60=\qty{60}{J},\qquad \Delta T=\frac{60}{4.18\times10^3\times60}=\qty{0.00024}{K} \]
\phtashikame
単位：Gy × kg＝J、J ÷ (J/(kg·K) × kg)＝K。\\
大きさ：60 J は、60 W の電球が1秒に出すエネルギーと同じで、熱としてはごくわずかです。\\
極端な場合：それでも、1 Gy を短い時間に全身に受けると、骨髄のしきい線量（約 0.5 Gy）をこえます。放射線の影響は熱ではなく、電離によるものです（④）。
\kotae \qty{60}{J}、約 \qty{0.00024}{K}
\end{phreidai}
\begin{jibun}
水 1 kg が 10 Gy を受けると、水の温度は何 K 上がりますか。
\bkjans{$10\div(4.18\times10^3)=\qty{0.0024}{K}$}
\end{jibun}

\label{ph:原-14:ex3}
\begin{phreidai}[α線とγ線の等価線量]
ある組織が、α線で 1 mGy、別のときにγ線で 1 mGy の吸収線量を受けました。それぞれの等価線量を求めましょう。
\phbamen
同じ吸収線量でも、放射線の種類で影響の大きさがちがうことを、数で確かめる場面です。
\phmushi
α線とγ線のエネルギーのちがいは考えず、放射線加重係数の表（表\ref{ph:原-14:tab-w}）の値を使います。
\phzu
$w_{\mathrm R}$：α線 20、γ線 1。
\phriyu
等価線量は、吸収線量に放射線加重係数をかけたものです（⑤）。
\phtateru
\[ H_{\mathrm T}=w_{\mathrm R}D_{\mathrm T} \]
\phtoku
\[ \text{α線：}20\times1=\qty{20}{mSv},\qquad \text{γ線：}1\times1=\qty{1}{mSv} \]
\phtashikame
単位：1 × mGy＝mSv（Gy も Sv も J/kg だが、重みをかけた量を Sv とよぶ）。\\
大きさ：α線は、同じエネルギーでも20倍の影響を見こみます。\\
極端な場合：α線は体の外からなら皮ふの表面で止まるので（原-13 ⑤）、この線量を受けるのは、主に体の中に取りこんだときです（⑧）。
\kotae α線 \qty{20}{mSv}、γ線 \qty{1}{mSv}
\end{phreidai}
\begin{jibun}
ある組織が中性子線で 2 mGy を受けました。このエネルギーの中性子の放射線加重係数を 10 とすると、等価線量はいくらですか。
\bkjans{$10\times2=\qty{20}{mSv}$}
\end{jibun}

\label{ph:原-14:ex4}
\begin{phreidai}[甲状腺の等価線量から実効線量へ]
甲状腺だけが、等価線量 10 mSv を受けました。ほかの組織は受けていません。実効線量はいくらですか。甲状腺の組織加重係数は 0.04 です。
\phbamen
体の一部だけが放射線を受けたとき、全身への影響を表す実効線量に直す場面です。
\phmushi
甲状腺以外の組織の線量を0とします。
\phzu
$H_{\text{甲状腺}}=\qty{10}{mSv}$、$w_{\mathrm T}=0.04$。
\phriyu
実効線量は、組織ごとの $w_{\mathrm T}H_{\mathrm T}$ の合計です（⑤）。
\phtateru
\[ E=\textstyle\sum w_{\mathrm T}H_{\mathrm T}=0.04\times H_{\text{甲状腺}} \]
\phtoku
\[ E=0.04\times10=\qty{0.4}{mSv} \]
\phtashikame
単位：1 × mSv＝mSv。\\
大きさ：甲状腺の等価線量は 10 mSv ですが、全身にならすと、その 4 \% にあたります。\\
極端な場合：全身が同じ 10 mSv を受ければ、$w_{\mathrm T}$ の合計が1なので、実効線量も 10 mSv です。
\kotae \qty{0.4}{mSv}
\end{phreidai}
\begin{jibun}
肺が 5 mSv、胃が 2 mSv の等価線量を受けました（ほかは0）。実効線量はいくらですか。
\bkjans{$0.12\times5+0.12\times2=\qty{0.84}{mSv}$}
\end{jibun}

\label{ph:原-14:ex5}
\begin{phreidai}[点線源からの距離と線量率]
セシウム137 の 1 MBq の点線源から、0.5 m・1 m・2 m の所の線量率（1 cm 線量当量率）を求めましょう。線量率定数は $\Gamma=\qty{0.0927}{\micro Sv.m^2/(MBq.h)}$ です。
\phbamen
放射線の授業で使う小さな線源（先生があつかう）から、どれだけ離れればよいかを考える場面です。
\phmushi
線源は点とし、空気による減弱、床や壁で散乱した放射線は無視します（⑥のモデル化の札）。自然放射線の分は数えません。
\phzu
$A=\qty{1}{MBq}$、$r=0.5$、1、\qty{2}{m}。
\phriyu
点線源の線量率は、距離の2乗に反比例します（⑥）。
\phtateru
\[ \dot H=\frac{\Gamma A}{r^2} \]
\phtoku
\[ r=0.5\text{：}\frac{0.0927}{0.25}=\qty{0.37}{\micro Sv/h},\qquad r=1\text{：}\frac{0.0927}{1}=\qty{0.093}{\micro Sv/h} \]
\[ r=2\text{：}\frac{0.0927}{4}=\qty{0.023}{\micro Sv/h} \]
\phtashikame
単位：\si{\micro Sv}·m$^2$/(MBq·h) × MBq ÷ m$^2$＝\si{\micro Sv/h}。\\
大きさ：2 m 離れると、自然放射線の空間線量率（0.1 \si{\micro Sv/h} 程度）の数分の1です。\\
極端な場合：線源にぴったり近づくと、点線源の式は使えなくなります（例題11）。
\kotae 0.5 m：\qty{0.37}{\micro Sv/h}、1 m：\qty{0.093}{\micro Sv/h}、2 m：\qty{0.023}{\micro Sv/h}
\end{phreidai}
\begin{jibun}
例題5 の線源から 3 m の所の線量率を求めましょう。
\bkjans{$0.0927/9=\qty{0.010}{\micro Sv/h}$}
\end{jibun}

\label{ph:原-14:ex6}
\begin{phreidai}[時間・距離・遮へいを全部使う]
セシウム137 の 37 MBq の点線源から 2 m の所に、厚さ 1 cm の鉛の板ごしに、30分いました。受ける線量を求めましょう。鉛の減弱係数は $\mu=\qty{1.26}{cm^{-1}}$（原-13 例題2）、$\Gamma=0.0927$ です。遮へいなしで 1 m の所に1時間いた場合とも比べましょう。
\phbamen
線源を使う作業で、3原則（⑦）で線量を見積もる場面です。
\phmushi
ビルドアップ（原-13 ⑩）、空気による減弱を無視します。
\phzu
$A=37$、$r=2$、$t=\qty{0.5}{h}$、$x=\qty{1}{cm}$。
\phriyu
距離は逆2乗、時間は比例、遮へいは $e^{-\mu x}$ で、3つをかけ合わせます（⑦の式）。
\phtateru
\[ H=\frac{\Gamma A}{r^2}\,t\,e^{-\mu x} \]
\phtoku
\[ \frac{0.0927\times37}{2^2}=\qty{0.858}{\micro Sv/h},\qquad e^{-1.26\times1}=0.284 \]
\[ H=0.858\times0.5\times0.284=\qty{0.12}{\micro Sv} \]
遮へいなしで 1 m に1時間なら $0.0927\times37=\qty{3.4}{\micro Sv}$ で、約28倍です。
\phtashikame
単位：\si{\micro Sv/h} × h＝\si{\micro Sv}。\\
大きさ：距離で $\tfrac14$、時間で $\tfrac12$、鉛で約 0.28 で、合わせて約 $\tfrac1{28}$。\\
極端な場合：ビルドアップを入れると、鉛の後ろの線量はもう少し大きくなります。
\kotae 約 \qty{0.12}{\micro Sv}（遮へいなし・1 m・1時間の約 $\tfrac1{28}$）
\end{phreidai}
\begin{jibun}
例題6 で、鉛を 2 cm にすると、線量はいくらになりますか。
\bkjans{$0.858\times0.5\times e^{-2.52}=\qty{0.034}{\micro Sv}$}
\end{jibun}

\label{ph:原-14:ex7}
\begin{phreidai}[自然放射線を1時間あたりに直す]
日本人が自然から受ける実効線量は、平均で1年に約 \qty{2.1}{mSv} です。1時間あたりでは何 \si{\micro Sv} ですか。
\phbamen
ニュースの \si{\micro Sv/h} と、1年の mSv を比べられるようにする場面です。
\phmushi
1年を 365.25 日とし、1年じゅう同じ割合で受けるとします（実際は場所や食事で変わる）。
\phzu
$\qty{2.1}{mSv}=\qty{2100}{\micro Sv}$、1年 $=365.25\times24$ 時間。
\phriyu
単位をそろえて、わり算をします。
\phtateru
\[ \frac{\qty{2100}{\micro Sv}}{365.25\times24\ \text{h}} \]
\phtoku
\[ \frac{2100}{8766}=\qty{0.24}{\micro Sv/h} \]
\phtashikame
単位：\si{\micro Sv} ÷ h。\\
大きさ：このうち外から受ける分は約 \qty{0.07}{\micro Sv/h} で、残りはラドンや食べ物から体の中で受ける分です（⑩）。\\
極端な場合：世界の平均 2.4 mSv/年なら \qty{0.27}{\micro Sv/h}。
\kotae 約 \qty{0.24}{\micro Sv/h}
\end{phreidai}
\begin{jibun}
空間線量率が \qty{0.05}{\micro Sv/h} の場所に、毎日8時間、1年（365日）いると、外から受ける線量はいくらですか。
\bkjans{$0.05\times8\times365=\qty{146}{\micro Sv}$、約 \qty{0.15}{mSv}}
\end{jibun}

\label{ph:原-14:ex8}
\begin{phreidai}[食品からの内部被ばく]
1 kg あたり 100 Bq（一般食品の基準値と同じ）のセシウム137 をふくむ食品を、大人が 1 kg 食べました。預託実効線量はいくらですか。線量係数は \qty{0.013}{\micro Sv}/Bq です。
\phbamen
食品の放射性物質の量（Bq）を、線量（Sv）に直す場面です。
\phmushi
食品にふくまれるのはセシウム137 だけとし、調理で減ることは無視します。
\phzu
濃さ 100 Bq/kg、量 1 kg、線量係数 0.013。
\phriyu
取りこんだ Bq に線量係数をかけます（⑧）。
\phtateru
\[ E=(100\times1)\times0.013 \]
\phtoku
\[ E=100\times0.013=\qty{1.3}{\micro Sv} \]
\phtashikame
単位：Bq × \si{\micro Sv}/Bq＝\si{\micro Sv}。\\
大きさ：自然放射線の約 5 時間分（0.24 \si{\micro Sv/h} × 5 h）です。\\
極端な場合：同じ食品を1年に 100 kg 食べると 0.13 mSv で、基準値を決めたときの「1年に 1 mSv」より小さくなります。
\kotae 約 \qty{1.3}{\micro Sv}
\end{phreidai}
\begin{jibun}
1 kg あたり 10 Bq（飲料水の基準値）のセシウム137 をふくむ水を、1日 2 L（2 kg）ずつ1年飲み続けると、預託実効線量はいくらですか。
\bkjans{$10\times2\times365\times0.013=\qty{95}{\micro Sv}$（約 0.095 mSv）}
\end{jibun}

\label{ph:原-14:ex9}
\begin{phreidai}[トリチウムをふくむ水を飲み続けたら]
1 L あたり 1500 Bq のトリチウムをふくむ水を、毎日 2 L、1年（365日）飲み続けたとします。預託実効線量はいくらですか。トリチウム（水の形）の線量係数は \qty{0.000018}{\micro Sv}/Bq です。
\phbamen
トリチウムの量（Bq）と線量（Sv）の関係を、具体的な数で確かめる場面です（ここでの濃さと量は、計算のために置いた仮の値です）。
\phmushi
飲む水以外からのトリチウムは数えません。1年分の取りこみを合計して計算します。
\phzu
濃さ 1500 Bq/L、1日 2 L、365日、係数 0.000018。
\phriyu
取りこんだ Bq の合計に線量係数をかけます（⑧）。
\phtateru
\[ E=(1500\times2\times365)\times0.000018 \]
\phtoku
\[ 1500\times2\times365=\qty{1.10e6}{Bq} \]
\[ E=1.10\times10^6\times1.8\times10^{-5}=\qty{20}{\micro Sv}=\qty{0.020}{mSv} \]
\phtashikame
単位：Bq × \si{\micro Sv}/Bq＝\si{\micro Sv}。\\
大きさ：自然放射線（2.1 mSv/年）の約100分の1です。\\
極端な場合：同じ Bq でも、セシウム137 なら線量係数が約700倍なので、線量も約700倍になります。Bq だけでなく、核種ごとの係数が大切です。
\kotae 約 \qty{0.020}{mSv}（\qty{20}{\micro Sv}）
\end{phreidai}
\begin{jibun}
1500 Bq の放射性物質を飲みこんだとき、それがトリチウム（水）なら、セシウム137 なら、預託実効線量はそれぞれいくらですか。
\bkjans{トリチウム $1500\times0.000018=\qty{0.027}{\micro Sv}$、セシウム137 $1500\times0.013=\qty{20}{\micro Sv}$（19.5）}
\end{jibun}

\label{ph:原-14:ex10}
\begin{phreidai}[背景のばらつきと検出限界]
食品の検査に使う測定器のバックグラウンドは、平均 30 cpm（1分に30カウント）です。10分はかったときの、1分あたりの数のばらつき（標準偏差）と、3σ 法の検出限界（背景をこえる分）を求めましょう。
\phbamen
「不検出」と「検出」の境目を決める場面です（⑨）。
\phmushi
バックグラウンドの平均は一定で、数のばらつきはポアソン分布にしたがうとします。
\phzu
10分で数える数 $N=30\times10=300$。
\phriyu
数えた数のばらつきは $\sqrt N$（③）。1分あたりに直すには10でわります。
\phtateru
\[ \sigma_{\text{1分あたり}}=\frac{\sqrt{N}}{10},\qquad 3\sigma \]
\phtoku
\[ \sqrt{300}=17.3,\qquad \sigma=\frac{17.3}{10}=\qty{1.7}{cpm},\qquad 3\sigma=\qty{5.2}{cpm} \]
\phtashikame
単位：cpm（1分あたりのカウント）。\\
大きさ：試料をはかって、背景より 5.2 cpm 以上多ければ「検出」。1分だけなら $3\sqrt{30}=\qty{16}{cpm}$ で、10分はかると約3分の1になります。\\
極端な場合：40分はかれば、さらに半分の約 2.6 cpm です（時間4倍で $\tfrac12$）。
\kotae $\sigma=\qty{1.7}{cpm}$、$3\sigma=\qty{5.2}{cpm}$
\end{phreidai}
\begin{jibun}
バックグラウンドが平均 100 cpm の測定器で4分はかるとき、1分あたりの σ と 3σ はいくらですか。
\bkjans{$N=400$、$\sigma=20/4=\qty{5}{cpm}$、$3\sigma=\qty{15}{cpm}$}
\end{jibun}

\label{ph:原-14:ex11}
\begin{phreidai}[検出器に飛びこむ割合：立体角]
半径 2 cm の円い窓の検出器を、点線源から 10 cm と 5 cm の所に、窓を線源に向けて置きます。線源から出る放射線のうち、窓に飛びこむ割合を、立体角の式と、逆2乗の近似 $a^2/(4d^2)$ で求めて比べましょう。
\phbamen
線源の放射能を、検出器の数から求めるときに要る「幾何学の効率」を見積もる場面です。
\phmushi
線源は点とし、空気や窓での減弱を無視します。
\phzu
$a=\qty{2}{cm}$、$d=10$、\qty{5}{cm}。
\phriyu
窓が線源から見て占める立体角の、$4\pi$ に対する割合です（⑥、8-8⑬）。
\phtateru
\[ \frac{\Omega}{4\pi}=\frac12\Bigl(1-\frac{d}{\sqrt{d^2+a^2}}\Bigr),\qquad \text{近似：}\frac{a^2}{4d^2} \]
\phtoku
\[ d=10\text{：}\frac12\Bigl(1-\frac{10}{\sqrt{104}}\Bigr)=0.0097\quad(\text{近似 }\tfrac{4}{400}=0.010) \]
\[ d=5\text{：}\frac12\Bigl(1-\frac{5}{\sqrt{29}}\Bigr)=0.036\quad(\text{近似 }\tfrac{4}{100}=0.040) \]
\phtashikame
単位：割合なので単位はありません。\\
大きさ：10 cm では 0.97 \% と 1.0 \% でほぼ同じ。5 cm では 3.6 \% と 4.0 \% で、近似が約1割大きくなります。\\
極端な場合：$d=2$ cm（窓の半径と同じ）では、正しい値 14.6 \% に対して近似は 25 \% で、まったく合いません。
\kotae 10 cm：0.97 \%（近似 1.0 \%）、5 cm：3.6 \%（近似 4.0 \%）
\end{phreidai}
\begin{jibun}
1秒に 1000 個のγ線を出す点線源から 10 cm の所に、例題11 の検出器を置きました。窓に飛びこむγ線は1秒に何個ですか。
\bkjans{$1000\times0.0097\approx9.7$ 個}
\end{jibun}

\label{ph:原-14:ex12}
\begin{phreidai}[線量率の高い場所で、何時間働けるか]
放射線業務従事者の線量限度は、5年で 100 mSv かつ1年で 50 mSv です。1年に平均 20 mSv までにおさめる計画で、線量率 \qty{10}{\micro Sv/h} の場所で作業するとき、1年に何時間まで作業できますか。
\phbamen
放射線を使う工場や発電所で、作業の計画を立てる場面です（⑫）。
\phmushi
ほかの場所での線量と、内部被ばくを0とします。線量率は一定とします。
\phzu
$\qty{20}{mSv}=\qty{20000}{\micro Sv}$、線量率 10 \si{\micro Sv/h}。
\phriyu
線量 ＝ 線量率 × 時間（⑦）です。5年で 100 mSv をこえないように、1年平均 20 mSv にします。
\phtateru
\[ t=\frac{\qty{20000}{\micro Sv}}{\qty{10}{\micro Sv/h}} \]
\phtoku
\[ t=\qty{2000}{h} \]
\phtashikame
単位：\si{\micro Sv} ÷ (\si{\micro Sv/h})＝h。\\
大きさ：1日8時間で250日分で、1年の仕事の時間に近い量です。\\
極端な場合：線量率が10倍（\qty{100}{\micro Sv/h}）なら200時間。距離や遮へいで線量率を下げれば、時間を増やせます。限度は「ここまで受けてよい」値ではないので、実際にはさらに下げる努力をします。
\kotae \qty{2000}{h}
\end{phreidai}
\begin{jibun}
女性の放射線業務従事者の限度（3か月で 5 mSv）では、\qty{10}{\micro Sv/h} の場所で3か月に何時間まで作業できますか。
\bkjans{$5000\div10=\qty{500}{h}$}
\end{jibun}

\label{ph:原-14:ex13}
\begin{phreidai}[100 mSv あたり 0.5 \% を、1000人に当てはめる]
1000人の集団の全員が 100 mSv を受けたとします。ICRP の防護の考え方（100 mSv あたり、がんで亡くなる確率が 0.5 \% 増える）を当てはめると、がんで亡くなる人は何人増えると計算されますか。もともと約30 \% の人ががんで亡くなることと比べましょう。
\phbamen
線量とリスクの関係を、集団の数で考える場面です（⑪）。
\phmushi
年齢や性別による感じやすさのちがいを無視し、集団の平均の値を使います。
\phzu
人数 1000、確率の増え方 0.5 \%、もとの割合 30 \%。
\phriyu
防護で使う LNT モデルで、確率を人数にかけます。
\phtateru
\[ (\text{増える人数})=1000\times0.005 \]
\phtoku
\[ 1000\times0.005=5\ \text{人},\qquad \text{もとの}\ 1000\times0.30=300\ \text{人} \]
\phtashikame
単位：人。\\
大きさ：300人が305人になる計算です。もとの300人も年や地域で変わるので、この5人の差を実際の調査で見分けることは、とてもむずかしいと考えられています。\\
極端な場合：100 mSv より少ない線量にこの計算を当てはめてよいかは、まだ決着がついていません（⑪のまだわかっていないこと）。
\kotae 約5人（もともとの約300人に加わる計算）
\end{phreidai}
\begin{jibun}
10万人の集団の全員が 10 mSv を受けたとして、LNT モデルで計算すると、がんで亡くなる人は何人増える計算になりますか。この計算の使い方について、ICRP はどう述べていますか。
\bkjans{$100000\times0.005\times0.1=50$ 人。ICRP は、ごく小さい線量を多くの人が受けたとしてがんの数を計算することは、公衆の健康の計画には適切でないとしている（⑪）}
\end{jibun}

%% ======================== まちがえやすいところ ========================
\Needspace{11\baselineskip}\bksection{tsumazuki}{まちがえやすいところをチェック}
\begin{tsumazuki}{Bq と Sv は、同じ物差し？}
\bkng{100 Bq の食品を食べると、100 \si{\micro Sv} の被ばくになる。}
\bkok{Bq は「1秒にこわれる数」（出す側）、Sv は「体が受けた影響の目安」（受ける側）。Bq を Sv に直すには、核種ごとの線量係数をかける。}
100 Bq のセシウム137 を食べても 1.3 \si{\micro Sv} です（例題8）。同じ Bq でも、核種によって Sv は何百倍もちがいます（例題9）。
\end{tsumazuki}

\begin{tsumazuki}{Gy と Sv は、いつも同じ数？}
\bkng{1 Gy を受けたら、いつも 1 Sv。}
\bkok{γ線・β線なら $w_{\mathrm R}=1$ で同じ数だが、α線は20倍。体の一部だけなら、実効線量は $w_{\mathrm T}$ の分だけ小さくなる。}
単位の名前がちがうのは、「重みをかけたかどうか」を区別するためです（⑤）。
\end{tsumazuki}

\begin{tsumazuki}{線量限度より下なら安全、上なら危険？}
\bkng{1年に 1 mSv までは安全で、1 mSv をこえると危険。}
\bkok{線量限度は、管理のための基準で、安全と危険の境界線ではない。限度より下でも、できるだけ下げる努力をする。}
影響の起きやすさは、線量とともに少しずつ変わります（⑪⑫）。
\end{tsumazuki}

\begin{tsumazuki}{測定器の数字は、ぴったり正しい？}
\bkng{1回目 0.06、2回目 0.05 と出たので、線量が下がった。}
\bkok{数えた数には $\sqrt N$ のばらつきがある。何回かはかって平均し、ばらつきの大きさと比べて判断する。}
長くはかるほど、ばらつきの割合は小さくなります（③⑨）。
\end{tsumazuki}

\begin{tsumazuki}{距離を2倍にすると、半分になる？}
\bkng{線源から2倍離れると、線量率は半分になる。}
\bkok{点線源なら、距離の2乗に反比例するので $\tfrac14$。ただし、広い地面に散らばった線源では、そうならない。}
逆2乗の法則は、点線源のモデルの式です（⑥のモデル化の札）。
\end{tsumazuki}

\begin{column}{放射線の単位になった人たち}
放射能の単位ベクレル（Bq）は、1896年にウランの塩から目に見えない放射線が出ていることを見つけたフランスのベクレルの名前です。グレイ（Gy）は、放射線が物に渡すエネルギーのはかり方を研究したイギリスのグレイ、シーベルト（Sv）は、放射線から人を守る研究と国際的なルールづくりに力をつくしたスウェーデンのシーベルトの名前です。ベクレルは、マリー・キュリーとピエール・キュリーとともに、1903年のノーベル物理学賞を受けました。放射線の発見から、その量をはかり、人を守るしくみができるまでには、多くの人の長い研究がありました。
\end{column}

%% ======================== 工学ではこう使う ========================
\Needspace{10\baselineskip}\bksection{phgear}{工学ではこう使う}
\begin{phkougaku}[放射線の管理：はかって、減らして、記録する]
\textbf{はかる}：放射線を使う病院・工場・研究所・原子力発電所では、部屋の空間線量率を電離箱やシンチレーション検出器で連続してはかり、働く人は個人線量計をつけて、受けた線量を毎日・毎月記録します。管理区域の出口では、体や物の表面に放射性物質がついていないかを、GM 計数管などで確かめます。

\textbf{減らす}：3原則（⑦）で、作業の手順（時間）、道具（距離：遠くから操作する装置）、遮へい（鉛の板・コンクリートの壁・水のプール）を工夫します。線量限度より下でも、できるだけ下げることが求められます（防護の最適化）。

\textbf{知らせる}：国や自治体は、各地のモニタリングポストの空間線量率をはかって公開しています。食品は、検出限界を決めて検査し、基準値と比べます（⑧⑨）。

\begin{phketa}[線量の目安]
\phketaline{自然放射線（日本の平均）}{約 2.1 mSv/年（約 0.24 \si{\micro Sv/h}）}
\phketaline{医療（日本の平均）}{約 2.6 mSv/年}
\phketaline{一般公衆の線量限度（ICRP）}{1 mSv/年（医療・自然は除く）}
\phketaline{放射線業務従事者}{50 mSv/年、100 mSv/5年}
\phketaline{確率的影響が見分けられる線量}{100 mSv ほどから上}
\phboxhead{答えが桁で合っているか}
\si{\micro Sv}・mSv・Sv は、それぞれ1000倍ずつちがいます。1時間あたり（/h）と1年あたり（/年）は、約 \num{8800} 倍ちがいます（1年 ≈ 8766 時間）。0.1 \si{\micro Sv/h} は、1年いると約 0.9 mSv です。答えが自然放射線と何桁ちがうかを、いつも確かめましょう。
\end{phketa}
\end{phkougaku}
\phstation{放射線管理,保健物理,医学物理,非破壊検査}{吸収線量、等価線量、実効線量、預託実効線量、1 cm 線量当量、管理区域、検出限界、不感時間、放射線透過試験}

%% ======================== 章末問題 ========================
\bkendproblems
\bkproblemgroup{基本}{3}
\bkproblem セシウム137 の 1 MBq の点線源から 1 m の所の線量率は約 \qty{0.093}{\micro Sv/h} です。4 m の所では、いくらになりますか。
\bkpans{距離4倍で $\tfrac1{16}$：$0.093\div16=\qty{0.0058}{\micro Sv/h}$}
\bkproblem 空間線量率が \qty{0.12}{\micro Sv/h} の場所で、1日8時間、30日すごすと、外から受ける線量は何 \si{\micro Sv} ですか。
\bkpans{$0.12\times8\times30=\qty{29}{\micro Sv}$（28.8）}
\bkproblem ある組織が、β線で 3 mGy、α線で 0.5 mGy の吸収線量を受けました。等価線量の合計を求めましょう。
\bkpans{$1\times3+20\times0.5=\qty{13}{mSv}$}
\bkproblem 1 kg あたり 50 Bq のセシウム137 をふくむ牛乳を、1日 1 kg ずつ30日飲んだときの預託実効線量を求めましょう（線量係数 \qty{0.013}{\micro Sv}/Bq）。
\bkpans{$50\times1\times30\times0.013=\qty{20}{\micro Sv}$（19.5）}
\bkproblem セシウム137 の 10 MBq の点線源から 1 m の所に1時間いると、受ける線量は約 \qty{0.93}{\micro Sv} です。距離を 2 m にし、時間を30分にすると、いくらになりますか。
\bkpans{距離で $\tfrac14$、時間で $\tfrac12$ なので $\tfrac18$：約 \qty{0.12}{\micro Sv}}
\bkproblem セシウム137 のγ線に対する鉛の半価層は 0.55 cm です。例題6 と同じ場面で、鉛を 1.1 cm にすると、鉛がないときの何分の1になりますか（ビルドアップは無視）。
\bkpans{半価層2枚分で $\tfrac14$}
\bkproblemgroup{チャレンジ}{4}
\bkproblem GM 計数管で1分はかったら \num{60000} カウントでした。不感時間を \qty{100}{\micro s} とすると、本当の数は1秒に何個ですか。数え落としは何 \% ですか。
\bkpans{$m=1000$ /s、$n=1000/(1-0.10)=\num{1111}$ /s（約 $\qty{1.1e3}{s^{-1}}$）、数え落とし 10 \%}
\bkproblem バックグラウンドが平均 20 cpm の測定器があります。（1）5分はかるとき、1分あたりの σ と 3σ を求めましょう。（2）3σ を 3 cpm 以下にするには、何分以上はかればよいですか。
\bkpans{（1）$N=100$、$\sigma=10/5=\qty{2}{cpm}$、$3\sigma=\qty{6}{cpm}$。（2）$3\sqrt{20t}/t\le3$ より $t\ge20$ 分}
\bkproblem 半径 1 cm の窓の検出器を、点線源から 1 cm の所に置きました。窓に飛びこむ割合を、立体角の式と、逆2乗の近似 $a^2/(4d^2)$ で求めて比べましょう。
\bkpans{立体角：$\tfrac12(1-1/\sqrt2)=0.146$（14.6 \%）。近似：$\tfrac14$（25 \%）。近すぎて近似は使えない}
\bkproblem 3 MBq のセシウム137 の点線源から 1 m の所に、鉛 1.1 cm（半価層2枚分）ごしに2時間いると、受ける線量はいくらですか（$\Gamma=0.0927$、ビルドアップは無視）。
\bkpans{$0.0927\times3\times2\times\tfrac14=\qty{0.14}{\micro Sv}$}

\begin{furikaeri}
\begin{checklist}
\item 放射線は目に見えないが、電離を使って数えられることと、Bq・Gy・Sv の意味のちがいを説明できる。\bkfurref{ph:原-14:sec2}
\item 吸収線量から等価線量・実効線量を計算し、自然放射線や医療の線量と比べられる。\bkfurref{ph:原-14:sec5}
\item 点線源の逆2乗の法則と遮へいを組み合わせ、時間・距離・遮へいの3原則で線量を見積もれる。\bkfurref{ph:原-14:sec7}
\item 計数のばらつき（$\sqrt N$）・数え落とし・検出限界・立体角を使って、測定の結果を正しく読める。\bkfurref{ph:原-14:sec3}
\item 確定的影響と確率的影響のちがい、線量限度の意味、低い線量のリスクについての2つの見方を説明できる。\bkfurref{ph:原-14:sec11}
\end{checklist}
\end{furikaeri}
\bkfukidashi{mina}{見えない放射線も、電離を使えば1個ずつ数えられるんだね。「危ないか安全か」じゃなくて、「どれだけの量か」を自然放射線と比べて考えるのが大事なんだ。}
